% 本文件由 LaTeX 排版 Agent 根据有效 translation.json 自动生成。
% author=Methodius; work_id=0623; title=十童女的宴席

\chapter{十童女的宴席}

% structure: p0001 source=h1 target=omitted reason=page-title-already-rendered-by-parent

\Emph{引言}\newline{}\Emph{讲道一}\newline{}\Emph{讲道二}\newline{}\Emph{讲道三}\newline{}\Emph{讲道四}\newline{}\Emph{讲道五}\newline{}\Emph{讲道六}\newline{}\Emph{讲道七}\newline{}\Emph{讲道八}\newline{}\Emph{讲道九}\newline{}\Emph{讲道十}\newline{}\Emph{讲道十一}\par

\SourceRecord{Translated by William R. Clark.；Ante-Nicene Fathers；Vol. 6.；Edited by Alexander Roberts, James Donaldson, and A. Cleveland Coxe.；Buffalo, NY: Christian Literature Publishing Co.,；1886.；<http://www.newadvent.org/fathers/0623.htm>.}

% structure: p0001 source=h1 target=section reason=page-title

\section{十童女的婚宴（引言）}

或论贞洁\par

对话人物：\Person{Euboulios}、\Person{Gregorion}、\Person{Arete}；\Person{Marcella}、\Person{Theophila}、\Person{Thaleia}、\Person{Theopatra}、\Person{Thallousa}、\Person{Agathe}、\Person{Procilla}、\Person{Thekla}、\Person{Tusiane}、\Person{Domnina}。\par

% structure: p0004 source=h2 target=subsection reason=relative-heading-rank; base=h2

\subsection{引言——作品计划；通往天堂之路；美德的描述与拟人化；\OriginalTerm{阿格诺斯}{Agnos}——贞洁的象征；\Person{Marcella}，基督童贞女中最年长、最卓越者。}

\Emph{Euboulios}：你来得正是时候，\Person{Gregorion}，因为我正找你，想听一听\Person{Marcella}和\Person{Theopatra}的会面，以及参加宴席的其他童贞女的情况，还有她们关于贞洁的谈话内容；因为据说她们的辩论如此精妙而有力，使得这个主题得到了充分的探讨。如果你来这里是为了别的事，那就把它推迟到另一个时间，不要耽搁，请给我们完整而连贯地讲述我们所询问的事情。\par

\Emph{Gregorion}：我似乎失去了希望，因为有人已经事先告诉了你我所询问的事情。我以为你对所发生的事一无所知，还沾沾自喜，以为我会是第一个告诉你的人。正因如此，我才急忙赶到这里来见你，担心的事果然发生了，有人抢先了一步。\par

\Emph{Euboulios}：放心，我亲爱的朋友，因为我们并没有得到任何确切的消息；那个带来消息的人什么也没告诉我们，只说有过交谈；但问他谈了什么，有什么用，他却不知道。\par

\Emph{Gregorion}：那么，既然我来这里正是为此，你是想从头听全部内容，还是让我略过一部分，只回忆那些我认为值得一提的要点？\par

\Emph{Euboulios}：千万不要略过；而是，\Person{Gregorion}，首先从最开始告诉我们聚会在哪里，酒席的摆设，以及你自己如何斟酒。\Quote{他们用金杯彼此祝饮，同时仰望广阔的天空。}\par

\Emph{Gregorion}：你总是善于辩论，在论证中极为有力——彻底驳倒所有对手。\par

\Emph{Euboulios}：现在不值得为这些事争执，\Person{Gregorion}；请你就简单地从头告诉我们发生了什么事。\par

\Emph{Gregorion}：好的，我试试。但先回答我这个问题：你大概认识\Person{Arete}，\Person{Philosophia}的女儿吧？\par

\Emph{Euboulios}：为什么这样问？\par

\Emph{Gregorion}：\Quote{我们受她邀请，前往她一个朝东的花园，享用当季的水果——我、\Person{Procilla}和\Person{Tusiane}。}我这是复述\Person{Theopatra}的话，因为我是从她那里得到消息的。\Quote{我们去了，\Person{Gregorion}，走了一条极为崎岖、陡峭而艰难的路；当我们靠近那地方时，} \Person{Theopatra}说，\Quote{迎面走来一位高大美丽的女子，从容优雅地走着，身穿一件雪白闪亮的长袍。她的美简直难以言表，超凡脱俗。她脸上绽放着谦逊与威严的融合。那是一张，}她说，\Quote{我从未见过的面孔，令人敬畏，却又透着温和与喜悦；因为它完全不加修饰，没有半点虚假。她向我们走来，像一个母亲见到久别的女儿，满心欢喜地拥抱并亲吻了我们每个人，说道：\InnerQuote{哦，我的女儿们，你们历经艰辛痛苦来到了我面前，我正热切地渴望引领你们进入不朽的牧场；你们艰难地走过一条布满许多可怕爬虫的道路；因为我曾看见你们常常偏离道路，害怕你们转身跌下悬崖。但感谢新郎，我已将你们许配给他，\BibleRef{格后 11:2}我的孩子们，因为他已让我们所有的祈祷得以有效成就。}正说着，} \Person{Theopatra}说，\Quote{我们来到了园子的围墙边，门还没有关，进去时，我们遇见了\Person{Thekla}、\Person{Agathe}和\Person{Marcella}正在预备晚餐。\Person{Arete}立刻说：\InnerQuote{你们也过来，和这些同伴一起坐在这里。}}现在，她对我说，\Quote{在那里做客的，我想一共有十位；那地方美丽非凡，充满娱乐设施。空气柔和而有规律地流动，混合着纯净的光线，一条如油般缓缓流动的溪流从花园正中央穿过，涌出最甘美的饮料；溪水清澈纯净，形成许多泉源，如河流般满溢，浇灌了整个花园；园中有各种树木，挂满新鲜果子，枝头果实盈盈，同样美丽；还有常青的草地，点缀着五彩缤纷、芳香四溢的花朵，一阵柔风拂过，带来最甜美的香气。\OriginalTerm{阿格诺斯}{Agnos}树长在旁边，一棵高大的树，因其极其宽阔而阴凉，我们在树下休憩。}\par

\Emph{Euboulios}：我的朋友，你似乎是在向我揭示第二个乐园。\par

\Emph{Gregorion}：你说得真实而明智。\Quote{到了那里，}她说，\Quote{我们有各种食物和多样的庆典，没有任何乐趣是缺乏的。之后，\Person{Arete}走进去，说了这些话：——}\par

\Quote{年轻的童贞女啊，我伟大荣耀的少女，美丽的童贞女，你们以未婚之手照料着基督无玷的牧场。我们现在已饱享饮食，因我们一切丰盛有余。那么，除了这个我还期望什么呢？我希望你们每一位都发表一篇赞颂童贞的讲论。让玛尔塞拉开始，因为她坐在首位，同时也是最年长的。若我不使成功的讲论者成为众人羡慕的对象，用永不凋谢的智慧之花为她加冕，我将为自己感到羞愧。}\par

\Quote{然后，}我想她说道，\Quote{玛尔塞拉立即开始发言如下。}\par

\SourceRecord{Translated by William R. Clark.；Ante-Nicene Fathers；Vol. 6.；Edited by Alexander Roberts, James Donaldson, and A. Cleveland Coxe.；Buffalo, NY: Christian Literature Publishing Co.,；1886.；<http://www.newadvent.org/fathers/062300.htm>.}

% structure: p0001 source=h1 target=section reason=page-title

\section{十童女的婚宴（首篇讲道）}

\Signature{玛尔塞拉}\par

% structure: p0003 source=h2 target=subsection reason=relative-heading-rank; base=h2

\subsection{第一章：童贞的艰难与崇高；论童贞女当研习教义。}

童贞是一件超性伟大、奇妙而荣耀的事；并且，平实地按照圣经来说，这最好最高贵的生活方式唯独是不朽之根，亦是其花与初果；为此，主在福音中论及人们自阉的种种原因时，应许那些自阉的人将进入天国\BibleRef{玛 9:12}。在人中，贞洁极为罕见且难以达到，而按其至高卓越与辉煌，其危险也极大。\par

因此，它需要刚强而宽宏的天性，这样的人跃过享乐之流，将灵魂的战车从地面向上驱驰，不偏离目标，直到因思想的迅捷轻快地超越世界，真正立于天穹之上，纯洁地默观从全能者无玷怀中涌出的不朽本身。\par

大地不能产生这饮品；惟独上天知道它所源出的泉源；因为我们应将童贞视为行在地上，却也伸及天上。为此，有些渴望童贞的人，只考虑到其结局，却因心志粗鄙，如同未洗脚一样徒劳而来，偏离了正道，因为他们对\Emph{童贞的}生活方式从未有过相称的观念。因为仅保持身体无玷是不够的，正如我们不应表现得更重视圣殿而非神像；我们却应照顾人的灵魂，视之为他们身体的神明，并以正义装饰它们。当人们不懈地努力聆听神圣的言论，直至磨光了智慧人的门槛\BibleRef{德 6:36}，达到认识真理时，他们才最是关心和照料它们。\par

因为正如肉体的腐液和脓液，以及所有那些使其腐败之物，都由盐驱除，同样，童贞女的一切非理性欲望也由神圣教导从身体中驱逐。因为未用基督的言语如盐般洒过的灵魂，必然发臭生虫，正如达味王在山中流泪公开承认时所喊的：\Quote{我的疮痍溃烂流脓}，因为他没有用克己的操练来腌自己，以克制肉欲，反而放任自己，在奸淫中腐败。因此在肋未纪中\BibleRef{肋 2:13}，凡不调盐的礼物，禁止作为祭品献给上主天主。如今圣经的一切属灵默想赐予我们作为盐，刺痛以得益，消毒以洁净；没有它，灵魂不可能藉理性被带到全能者面前；因为主对宗徒说：\Quote{你们是地上的盐}\BibleRef{玛 5:13}。因此，童贞女应当始终爱慕尊贵之事，在智慧上出类拔萃，不沉溺于任何懒惰或奢侈之事，而应卓越超群，专心于与童贞身份相称的事，常常藉言语除去奢侈的污秽，免得些许隐藏的腐败滋生不洁之虫；因为蒙福的保禄说：\Quote{未出嫁的女子挂虑主的事，为叫身心圣洁}\BibleRef{格前 7:34}。但她们中许多人视听道为次要之事，若稍加留意，便以为行了大事。然而对这些事必须加以辨别；因为将神圣教训传授给一个注重琐碎、卑下并假冒智慧的天性，是不合适的。因为对那些将全部精力用于小事，以便把他们想要完善的东西精确完成，却不认为最需要操劳的是那些最能增加贞洁之爱的必要之事的人，继续与他们谈话，岂不可笑？\par

% structure: p0008 source=h2 target=subsection reason=relative-heading-rank; base=h2

\subsection{第二章：童贞乃天上之植物，晚近始引入；人类迈向成全的进步，如何安排。}

的确，借着极大的力量，童贞的植物从天降下给人，正因如此，它并未启示给最初几代人。因为人类种族当时数量极少；必须先让数目增加，然后才能达至完美。因此，古人并不认为娶自己的妹妹为妻有何不当，直到法律来到，将他们分开，并禁止那起初看来正确的事，宣布其为罪恶，称那\Quote{揭露他姊妹的裸体}的人为可咒骂的；天主这样仁慈地按季节将必要的帮助带给我们的种族，如同父母对待他们的孩子。因为父母不会立即为孩子设立师长，而是在孩童时期允许他们像小动物一样玩耍，先送他们到与自己一样咿呀学语的教师那里，直到他们抛弃心智上的稚嫩绒毛，进而实践更大的事，再从那里实践更大的事。同样，我们必须认为，万有的天主和圣父对待我们的先祖也是如此。因为世界在尚未充满人类时，就像一个孩子，必须先被人类充满，然后才长大成人。但此后，当世界从一端到另一端都被殖民，人类种族无限扩展时，天主不再允许人停留在同一种方式中，他考虑他们如何能从一点进到另一点，更接近天堂，直到他们获得最伟大、最崇高的童贞教训，达至完美：首先，他们应放弃兄妹通婚，而从其他家族娶妻；然后，不应再像禽兽那样拥有多妻，仿佛只为繁殖而生；然后，不应有奸淫；再然后，他们应继续节制，从节制到童贞，当他们训练自己轻视肉体时，便无畏地驶入永生的平静港湾。\par

% structure: p0010 source=h2 target=subsection reason=relative-heading-rank; base=h2

\subsection{第三章 因\Person{亚巴郎}的割损，禁止与姊妹通婚；在先知时代，多妻制被停止；婚姻的纯洁本身逐步被强制。}

如果任何人敢于指责我们的论证缺乏圣经证据，我们将提出先知著作，更充分地证明已做的陈述的真实性。当\Person{亚巴郎}初次接受割损的盟约时，他似乎通过在自己身体的一个成员上接受割损，所表明的无非是：人不应再与同父母所生的人生育子女；表明每个人都应避免与他自己的姊妹同房，如同避免自己的血肉一样。因此，从\Person{亚巴郎}的时代起，与姊妹通婚的习俗已经终止；从先知的时代起，与多个妻子缔结婚姻已被废除；因为我们读到：\BibleQuote{不要随从你的私欲，但要抑制你的贪欲；} \BibleRef{德 18:30} 因为\BibleQuote{酒色能使明智人堕落；} \BibleRef{德 19:2} 在另一处：\BibleQuote{让你的泉源蒙受祝福；并要与你年轻时所娶的妻子欢乐，} \BibleRef{箴 5:18} 明显禁止多妻。\Person{耶肋米亚}清楚地将那些贪恋其他妇女的人称为\BibleQuote{喂饱的马} \BibleRef{耶 5:8}；我们又读到：\BibleQuote{不虔诚者的众多后代必不繁荣，从杂种枝条也不能生根深厚，也不能奠定任何牢固的基础。} \BibleRef{智 4:3}\par

然而，免得我们因收集先知们的见证而显得冗长，让我们再次指出，贞洁是如何继一妻制之后出现的，逐步除去肉体的情欲，直至完全消除因习惯而产生的性交倾向。因为不久，有人被介绍进来，恳切地祈求从此远离这种诱惑，说：\Quote{上主，我生命之父和治理者，不要让我顺从他们的计谋；不要赐给我傲慢的眼神；不要让肚腹的贪婪和肉体的情欲控制我。} 在充满一切德行的智慧篇中，圣神如今公开吸引听众趋向节制和贞洁，如此歌唱：\BibleQuote{没有子女而有德行是更好的，因为德行的纪念是不朽的；因为它为天主及人所共知。当德行在场时，人们受它感动；当它离去时，人们渴望它：它永远戴着冠冕，得胜凯旋，赢得了无玷的奖赏。} \BibleRef{智 4:1}\par

% structure: p0013 source=h2 target=subsection reason=relative-heading-rank; base=h2

\subsection{第四章 惟独\Person{基督}教导了童贞，公开宣讲天国；应在神圣德行之光中达到天主的肖像。}

我们已论及人类的各种时期，并如何它始于兄弟姊妹的通婚，进而达到节制；现在留给我们的是童贞的主题。那么，让我们尽我们所能来谈论它。首先，让我们探究为何许多族长、先知和义人，他们教导并做了许多崇高的事，却无人称赞或选择童贞的状态。因为童贞是保留给主独自首先教导此道理的，唯独祂降来人间，教导人亲近天主；因为那身为司祭、先知和天使之首与元首者，也应当被尊为童贞之首与元首。因为在古时，人尚未完美，因此不能接受那完美的，即童贞。因为人受造是天主的\Emph{肖像}，还需要获得那相似祂的\Emph{相似}；为此圣言被派遣到世界上来使此完美。祂首先取了我们的形状，这形状被许多罪恶所扭曲，好使我们因祂的缘故——祂承担了它——得以再次接受天主的\Emph{形象}。因为当我们如同熟练的画师，将祂的特征印在我们身上如同印在版上，学习祂为我们所指示的道路时，我们才真正是按照天主的相似而被塑造。为此，祂身为天主，乐意穿上人的血肉，好使我们如同在版上看见我们生活的神圣典范，也能模仿那描绘它的祂。因为祂并非心口不一之人；也不是认为某事正确却教导另一事。凡是真正有益和正确的事，祂都既教导又实行。\par

% structure: p0015 source=h2 target=subsection reason=relative-heading-rank; base=h2

\subsection{第五章 基督，通过保守祂的肉体在童贞中不受腐坏，吸引人修习童贞；童贞者的人数与圣人人数相比的比例}

那么，那位是真理和光明的主，从天上降下来时着手做了什么呢？祂保守了祂所取来的肉体在童贞中不受腐坏，好使我们若愿意相似天主和基督，也应努力尊崇童贞。因为天主的相似就是避免腐坏。那拥有基督的\Person{若望}向我们展示，在\BibleBook{}{默示录}中说：\BibleQuote{我观看，看，羔羊站在熙雍山上，同祂在一起的十四万四千人，都有祂的名和祂父的名写在额上。我听见一个声音从天上传来，好像众水的声音，又像大雷的声音；我所听见的声音，好像弹竖琴者弹着他们的竖琴。他们在宝座前、四个活物和长老面前唱新歌；除了那从地上被赎的十四万四千人，没有人能学会那歌。这些人是未曾与女人玷污的，因为他们乃是童贞者。这些人是跟随羔羊的，无论祂到哪里，他们都跟随。}\BibleRef{默 14:1}显示主是童贞者歌咏团的领唱者。此外，请注意，在天主眼中童贞的尊荣是多么伟大：\BibleQuote{这些人是从人类中被赎的，成为天主和羔羊的初果。在他们口中找不出诡诈，因为他们是没有瑕疵的，}\BibleRef{默 14:4}他说，\BibleQuote{他们跟随羔羊，无论祂到哪里。}他显然想借此教导我们，从一开始，童贞者的数目就被限制为这么多，即十四万四千，而其他圣人的群众则是不可胜数的。因为让我们思考他在论及其余人时所说的话：\BibleQuote{我看见一大群人，没有人能数得过来，来自各民族、各支派、各人民、各语言。}\BibleRef{默 7:9}因此，正如我所说，显然在论及其他圣人时，他引入了一个不可言喻的群众，而在论及那些处于童贞状态的人时，他只提到一个非常小的数目，以便与那些构成无数数目的人形成强烈对比。\par

\Person{阿瑞忒}啊，这就是我关于童贞主题的论述。但如果我遗漏了什么，让继我之后的\Person{德奥斐拉}来补充遗漏。\par

\SourceRecord{Translated by William R. Clark.；Ante-Nicene Fathers；Vol. 6.；Edited by Alexander Roberts, James Donaldson, and A. Cleveland Coxe.；Buffalo, NY: Christian Literature Publishing Co.,；1886.；<http://www.newadvent.org/fathers/062301.htm>.}

% structure: p0001 source=h1 target=section reason=page-title

\section{\WorkTitle{十贞女的盛宴（第二讲）}}

\Person{Theophila}\par

% structure: p0003 source=h2 target=subsection reason=relative-heading-rank; base=h2

\subsection{\Emph{第一章　赞美贞洁并未废除婚姻。}}

于是她说道，\Person{Theophila}发言：既然\Person{Marcella}卓越地开始了这一讨论，却未能充分完成，那么我有必要尽力为其画上句号。人逐步迈向贞洁，天主不时地催促他，这一点在我看来已被绝妙地证明；但我不能对以下断言持同样看法：即从今以后他们不应再生育子女。因为我以为我从圣经中清楚地看到，在引入贞洁之后，\OriginalTerm{圣言}{the Word}并未完全废除生育子女；因为虽然月亮可能比星辰更大，但其他星辰的光辉并未被月光消灭。\par

让我们从\BibleBook{}{创世纪}开始，以便将古远和至高地位赋予这部经文。如今，关于生育子女的天主命令和规定\BibleRef{创 1:28}确实在今日仍得应验，造物主仍在塑造人。因为这一点十分明显：天主如同一位画师，此刻仍在世界工作，正如主所教导的：\Quote{我父至今工作。}但当江河止息、不再流入大海的蓄池，光明与黑暗完全分离——因为分离仍在进行——旱地不再生出爬行动物和四足兽类的果实，预定数目的人已满，那时人将不再生育子女。但如今，在世界存在且仍在形成之时，人必须配合天主塑造祂的肖像；因为经上说：\Quote{你们要生育繁殖。}\BibleRef{创 1:28}我们不可对造物主的命令感到反感，因为我们自己也正是由此而来。因为将种子播撒在子宫的犁沟中，是人生育的开端，如此，骨由骨、肉由肉，藉着一种无形的力量，被塑造成另一个人。这样，我们必须认为这句话得到了应验：\Quote{这才真是我的骨中之骨，肉中之肉。}\BibleRef{创 2:23}\par

% structure: p0006 source=h2 target=subsection reason=relative-heading-rank; base=h2

\subsection{\Emph{第二章　生育类似于从亚当的肋骨和本性形成厄娃；天主在通常生育中创造人。}}

这也许就是第一个人的沉睡和出神所预表的，它预示了婚姻爱情的拥抱。当人渴望子女时，他陷入一种出神，被生育的快乐软化并制服，如同沉睡，以至于从他的肉和骨中又取出一部分，如我所说，被塑造成另一个人。因为身体在爱情的拥抱中失去和谐，正如那些有婚姻经验的人告诉我们的，所有像骨髓一样具有生殖力的血从各个肢体汇集而来，形成泡沫并凝结，通过生育器官投射到女性的活体内。大概正是因为这个原因，人说男人要离开父亲和母亲，因为当他与妻子在爱情的拥抱中结合时，他突然忘记一切，被生育的欲望所征服，将自己的肋骨献给神圣的造物主，以便从中取走，使父亲在儿子中再次出现。\par

因此，如果天主仍在塑造人，如果我们认为生育子女是可憎的，而全能者自己却不羞于以祂无玷的手使用它，那么我们岂不犯有鲁莽之罪？因为祂对\Person{耶肋米亚}说：\Quote{我还没有在母腹中形成你以前，我已认识了你；}\BibleRef{耶 1:5}对\Person{约伯}说：\Quote{你岂取泥土捏造一个活物，使它在地上说话？}而\Person{约伯}近前恳求说：\Quote{你的手制造我、塑造我。}\BibleRef{约 10:8}那么，禁止婚姻结合岂不是荒谬的吗？因为我们期望在我们之后将有殉道者和对抗那恶者的人，为此\OriginalTerm{圣言}{the Word}也应许要缩短那些日子。\BibleRef{玛 24:22}因为如果生育子女从此在天主眼中被视为邪恶，如你所说，那么那些违背神圣法令和旨意而存在的人，又如何能够取悦天主呢？所生的岂非是私生子，而非天主的受造物，若像假币一样，在法律权威的意图和命令之外被塑造？因此，我们承认人拥有塑造人的能力。\par

% structure: p0009 source=h2 target=subsection reason=relative-heading-rank; base=h2

\subsection{\Emph{第三章　一段模糊的经文；不仅信徒，甚至主教有时也是不合法的。}}

但\Person{Marcella}打断说：啊，\Person{Theophila}，这里似乎有一个大错误，与你所说的相悖；你以为能躲在你抛出的云雾之下吗？因为来了那个论据，或许任何视你为极有智慧的人都会提出：你如何评价那些在通奸中非法所生的人？因为你认定，任何人除非由神圣统治者的旨意引入，其身体由天主预备，否则不可能进入世界，这是不可思议和不可能的。而为了不让你躲在安全的墙后，提出那节经文说：\Quote{至于奸夫所生的子女，他们不得达到完美，}\BibleRef{智 3:16}他轻易回答你说：我们常常看到那些非法所生的人如成熟的果实般达到完美。\par

如果，你又一次诡辩地回答：\Quote{哦，我的朋友，那些不达完美的人，我理解为是在基督教导的义德中得以成全的；}那么他会说：\Quote{但是，我尊贵的朋友，由不义的种子所生的人，不仅被列入聚集在弟兄羊群中的人中，而且常常甚至被召来管理他们。既然这一点是清楚的，且所有人都证明，由奸淫所生的人确实达至完美，我们就不应想象圣神是在教导有关怀孕和出生的事，而是更可能是在教导那些篡改真理的人，他们以虚假教义败坏圣经，产生不完美且不成熟的智慧，将他们的错误与虔诚相混合。}因此，这一辩解从你那里被除去后，现在来告诉我们，由奸淫所生的人是否由天主的旨意所生？因为你说过，除非主塑造并赐予生命，人的后代不可能达至完美。\par

% structure: p0012 source=h2 target=subsection reason=relative-heading-rank; base=h2

\subsection{第四章 人的生育，及其中天主的工作}

泰奥菲拉好像被一个强大的对手拦腰抱住，感到晕眩，好不容易镇定下来，回答说：\Quote{我的尊贵朋友，你提的问题需要用比喻来解决，好让你更能理解天主的创造能力如何遍及万物，尤其在人的生育中乃是真正的成因，使那些播种在丰饶大地上的作物生长。因为所播种的本身不应受责备，而是那以非法拥抱在陌生土壤中播种的人，仿佛以可耻的方式出卖自己的种子来换取些许快感。试想，我们出生到这世界，好比一所房子，其入口紧邻高山；这房子向下延伸很远，远离入口，后面有许多洞，而且这一部分是环形的。}玛尔塞拉说：\Quote{我想象得到。}\Quote{那么，再假设里面有一位塑造者，正在雕刻许多塑像；再想象，泥土的原料不断地从外面通过那些洞，由许多根本无法看见那艺术家本人的人送入。再假设这房子被云雾笼罩，外面的人什么也看不见，只有那些洞。}她说：\Quote{这也可以假设。}而且，那些一起提供泥土的人中，每人都有一个指定的洞，他只能把自己的泥土带进去并存放，不能碰其他任何洞。如果他又擅自试图打开分配给他人的洞，就要受到火与鞭打的威胁。\par

\Quote{好，现在再考虑接下来发生的事：里面的塑造者到每个洞去，私下取走他在每个洞找到的泥土用于塑造，经过一定月份塑成后，通过同一个洞交还；他遵循这条规则：每一块能被塑造的泥土都要不加区别地加工，即使有人非法地通过别人的洞扔进来，因为泥土没有犯错，因此作为无辜者，应当被塑造和成形；但那个违反规定和法律将其投放在别人洞中的人，应作为罪犯和违法者受罚。因为泥土不应受责备，而是那违反正义做这事的人；由于不节制，他偷走了种子，秘密地、暴力地投放到别人的洞中。}\Quote{你说得极对。}\par

% structure: p0015 source=h2 target=subsection reason=relative-heading-rank; base=h2

\subsection{第五章 圣教父继续同一论述}

既然这些事已经完成，我最聪明的朋友们，你们还需将这画面应用于已经谈过的事；将房子比作我们生育的不可见本性，将靠山的入口比作灵魂从天降下和进入身体的过程；洞比作女性，塑造者比作天主的创造能力，它在生育的掩护下，利用我们的本性，无形中在我们内塑造人，为灵魂制作外衣。搬运泥土的人比喻男性；当他们渴求孩子时，将种子带入女性的自然通道，正如比喻中那些人将泥土投入洞中。因为种子，可以说分享天主的创造能力，不应被视为不节制之冲动的罪魁祸首；技艺总是加工交给它的材料；没有什么是本身邪恶的，而是因使用者的行为而成为邪恶；因为当正当而纯洁地使用时，它变得纯洁，但若可耻而不正当地使用，就变得可耻。铁原为农业和技艺的益处而被发现，它如何伤害那些将其磨快用于杀戮战场的人？金、银、铜，以及总的说来所有可加工的泥土，又如何伤害那些忘恩负义对待造物主、错误地将它们的一部分转变成各种偶像的人？如果有人将从偷盗得来的羊毛供应给织布技艺，该技艺只关注一点，即加工所给的材料，只要它可接受加工，不拒绝任何对它有用的东西，因为被偷的羊毛本身没有生命，不应受责备。因此，材料本身应被加工和装饰，而发现不正当地偷取材料的人应受惩罚。同样，婚姻的破坏者、打破生命和谐之弦（如同竖琴之弦）的人，因情欲狂怒，在奸淫中放纵欲望，他们应受折磨和惩罚，因为他们从他人的花园中偷取生育的拥抱，犯了大错；但种子本身，如同羊毛，应被塑造并赋予生命。\par

% structure: p0017 source=h2 target=subsection reason=relative-heading-rank; base=h2

\subsection{第六章 天主甚至眷顾奸淫所生的子女；赐予他们天使作护守}

但何需援引这些例子来延长论证呢？因为若没有天主的帮助，自然不可能在短时间内完成如此伟大的工程。是谁赋予了骨骼固定的本性？是谁用神经约束了柔软的肢体，使其能在关节处伸展和收缩？是谁为血液准备了管道，为气息准备了柔软的气管？又是什么神祇促使体液发酵，将其与血液混合，并从泥土中形成柔软的肉体？岂不唯独是至高的工匠，祂使我们成为人，祂自己理性而活生生的肖像，如同用蜡在子宫中从潮湿微小的种子塑造而成？或凭谁的眷顾，胎儿在封闭于血管连接处时没有被潮湿窒息？或在它诞生并见到光明之后，是谁使其从柔弱微小变得高大、美丽而强壮？岂不正是天主自己，至高的工匠，如我所说，藉其创造的大能制作基督的复制品和生动的画像？因此我们也从受启示的著作中得知，那些被生下来的人，即使是在奸淫中所生，也被交托给守护天使。但如果他们的诞生违背了天主福乐本性的意愿和法令，他们又怎会被交付给天使，以极大的温柔和宽容被养育呢？又如果他们能控告自己的父母，怎能自信地在基督的审判座前呼求祂说：\Quote{主啊，祢没有吝惜这共同的光明给我们；但这些人藐视祢的命令，定我们于死亡。}因为\Quote{祂说：}\InnerQuote{非法床榻所生的子女，在受审时，是控告父母邪恶的证人。}\BibleRef{智 4:6}\par

% structure: p0019 source=h2 target=subsection reason=relative-heading-rank; base=h2

\subsection{第七章 理性灵魂来自天主本身；贞洁并非唯一的善，尽管是最佳且最受尊崇的。}

或许在缺乏辨别力和判断力的人中，有人可以看似有理地争辩说：这由人所栽种的灵魂之肉身外衣，是在没有天主判决的情况下自行成形的。然而，如果他还教导说：灵魂的不朽存有也与可朽的肉身一同被播下，那将无人相信；因为唯有全能者将不朽不腐的部分吹入人内，正如也只有祂是不可见和不可毁灭的创造主。因为祂说：祂\Quote{向他的鼻孔吹了一口生气，人就成了一个有灵的生物。}\BibleRef{创 2:7}那些制造人形偶像以毁灭人类的工匠，没有察觉也不认识他们自己的造主，被圣言所责备，在充满一切德行的\BibleBook{}{智慧篇}上说：\Quote{他的心是灰烬，他的希望比尘土更卑贱，他的生命比泥土更无价值；因为他认识了自己的造主，未曾认识那赋予他行动的灵魂，并注入生活精神的那一位，}\BibleRef{智 15:10}即天主，众人的创造主；因此，按照宗徒的话，祂\Quote{愿意所有人都得救，并得以认识真理。}\BibleRef{弟前 2:4}现在，虽然这个主题几乎尚未完成，但仍有一些需要讨论的。因为当人彻底考察并理解那些按照其本性发生在人身上的事时，他将不会轻视生育子女，尽管他赞扬贞洁并尊崇它。因为虽然蜂蜜比其他东西更甜、更令人愉快，但我们不能因此认为那些混合在果实天然甘甜中的其他东西是苦的。为了支持这些说法，我将提出一个可靠的见证人，即\Person{保禄}，他说：\Quote{这样看来，那结婚的，做得好；那不结婚的，做得更好。}\BibleRef{格前 7:38}此话在指出那更美好更甘甜之事时，并非意图除去较低者，而是安排给每一样事物各自适当的用途和益处。因为有些人未能获得童贞；另一些人祂不再愿意他们因生育而被激起情欲、被玷污，而从此要默想并专注于身体改变为天使的相似，那时他们\Quote{也不娶，也不嫁，}\BibleRef{玛 22:30}按照主可靠的言语；因为并非所有人都能达到为天国自阉的那种无玷状态，\BibleRef{玛 19:12}而显然只有那些能保持童贞这常青不谢的花朵的人。因为预言性的圣言惯于将教会比作一片鲜花覆盖、五彩缤纷的草地，不仅以童贞之花，也以生育和节制之花装饰和加冕；因为经上记载：\Quote{王后佩戴着金饰，身穿锦绣衣袍，站在你的右边。}\par

\Person{阿勒特}啊，这些言语是我为真理尽己所能在这场讨论中带来的。\par

当\Person{泰奥斐拉}这样说时，\Person{泰奥帕特拉}说，所有赞同她言论的童女中响起了掌声；她们安静下来后，经过长时间的停顿，\Person{塔雷娅}站起来，因为她在竞赛中被分配了第三名，紧随\Person{泰奥斐拉}之后。于是她，如我所想，接着发言了。\par

\SourceRecord{Translated by William R. Clark.；Ante-Nicene Fathers；Vol. 6.；Edited by Alexander Roberts, James Donaldson, and A. Cleveland Coxe.；Buffalo, NY: Christian Literature Publishing Co.,；1886.；<http://www.newadvent.org/fathers/062302.htm>.}

% structure: p0001 source=h1 target=section reason=page-title

\section{\WorkTitle{十童女的宴会（第三篇讲道）}}

\Person{塔莱娅}\par

% structure: p0003 source=h2 target=subsection reason=relative-heading-rank; base=h2

\subsection{第一章 \WorkTitle{圣经}章节的比较}

\Person{特奥菲拉}啊，在我看来，你在行动和言语上都超越众人，在智慧上也不逊于任何人。因为无论一个人多么好争论和矛盾，都不会指责你的讲论。然而，尽管其他一切都似乎讲得正确，但有一件事，我的朋友，使我感到痛苦和困扰，因为我考虑到那位明智而最属灵的人——我指的是\Person{保禄}——不会徒然地将第一个男人和女人的结合归于\Person{基督}和教会，如果圣经所表达的不比单纯的词语和历史所传达的更高；因为如果我们把圣经仅仅看作是完全涉及男女结合的表象，那么宗徒为何要回忆这些事，并引导我们（如我所认为）进入圣神的道路，将\Person{亚当}和\Person{厄娃}的历史寓意化，视为指向\Person{基督}和教会呢？因为\BibleBook{}{创世纪}中的经文是这样写的：\BibleQuote{亚当说：\InnerQuote{这是我骨中的骨，肉中的肉；可以称她为女人，因为她是从男人身上取出来的。因此，人要离开父母，与妻子连合，二人成为一体。}} \BibleRef{创 2:23} 但是宗徒考虑这段经文，绝不像我所说的那样，意图按其单纯的自然意义来理解它，作为男女结合，正如你所做的；因为你以过于自然的意义解释这段经文，断定圣神只是谈论怀孕和生育；从骨中取出的骨成了另一个人，活物在怀孕时像树一样膨胀。但他更属灵地将这段经文指向\Person{基督}，这样教导：\BibleQuote{爱妻子的就是爱自己。从来没有人恨恶自己的肉身，总是保养顾惜，正像主对待教会一样；因为我们是他身体的肢体，是他的肉，是他的骨。因此，人要离开父母，与妻子连合，二人成为一体。这是极大的奥秘；但我是指着基督和教会说的。}\par

% structure: p0005 source=h2 target=subsection reason=relative-heading-rank; base=h2

\subsection{第二章 宗徒\Person{保禄}的离题；其教义的特点：毫无矛盾；谴责\Person{奥利振}，他错误地将一切都转为寓意}

如果他在讨论一类主题时，转入另一类，以致似乎混淆它们，并引入与所讨论主题无关的事物，偏离问题，正如现在这样，请不要感到困扰。因为似乎为了最谨慎地加强关于贞洁的论证，他事先准备了论证方式，从更具说服力的讲话开始。因为他讲话的风格非常多样，且为逐步证明而安排，开始温和，但流向更高贵、更宏伟的风格。然后，又转向深奥之处，有时以简单易懂的方式结束，有时以更困难、更精细的方式结束；然而通过这些变化，他没有引入任何与主题无关的东西，而是根据某种奇妙的关系将它们全部聚集在一起，将提出的问题融为一炉。因此，我需要更准确地解释宗徒论证的含义，同时不拒绝之前所说的一切。因为在我看来，\Person{特奥菲拉}啊，你似乎已经充分而清晰地讨论了圣经中的那些话，并且无误地陈述了它们。因为完全轻视字面意义是危险的，尤其是\BibleBook{}{创世纪}，其中天主为宇宙的构成所定的不变旨意被展示出来，与此一致，直到如今，世界都完美地有序，根据完美的规则，极其美丽地运行，直到立法者亲自重新安排，希望重新整顿，通过新的安排打破自然界最初的法则。但是，既然不宜让论证的证明未经审查——可以说是半跛的——来吧，让我们，如同完成我们的配对，提出类比的意义，更深入地审视圣经；因为当\Person{保禄}越过字面意义，显示话语延伸至\Person{基督}和教会时，他是不应被轻视的。\par

% structure: p0007 source=h2 target=subsection reason=relative-heading-rank; base=h2

\subsection{第三章 第一个亚当与第二个亚当的比较}

首先，我们必须探究\Person{亚当}是否能比作\Person{天主子}，当他处于堕落的过犯中，并听到判决：\BibleQuote{你本是尘土，仍要归于尘土。} \BibleRef{创 3:19} 因为，他怎么会被视为\BibleQuote{首生的，在一切被造之先} \BibleRef{哥 1:15}，那在天地之后用泥土所造的呢？他又如何能被承认为\BibleQuote{生命树}，因过犯被赶出，\BibleRef{默 2:7} 免得\BibleQuote{他再伸手摘取，永远活着}？\BibleRef{创 3:22} 因为与某物相似的事物，必须在许多方面与其相似物类似并相称，而不能有相反和不相似的性质。因为若有人胆敢将不匀的比作匀的，或和谐比作不和谐，则不会被认为是理性的。但匀的应与本质上匀的相比，即使只是稍微匀；白应与本质上白的相比，即使非常小，并且仅适度显示出使其被称为白的白度。现在，毫无疑问，每个人都很清楚，无罪的和不朽的是匀称、和谐、明亮如智慧；但可朽的和有罪的是不匀称、不和谐，被驱逐出去，有罪且应受谴责。\par

% structure: p0009 source=h2 target=subsection reason=relative-heading-rank; base=h2

\subsection{第四章 某些难题以及处理过简，似乎未按神学规则充分阐明的部分}

因此，我认为这就是许多人提出的反对意见，他们似乎轻视\Person{保禄}的智慧，不喜欢将第一个人与\Person{基督}相比较。来吧，让我们思考\Person{保禄}如何正确地比较了\Person{亚当}和\Person{基督}，不仅视他为预像和肖像，而且\Person{基督}自己成了完全相同的事，因为永恒圣言降临在他身上。因为理当如此：天主的首生子、初芽、独生子，甚至天主的智慧，应与最初受造的人、人类的首生者联合，并降生成人。这就是\Person{基督}，一个人充满纯洁而完美的天主性，而天主被接纳到人内。因为最合适的是，最古老的永世和首位总领天使，在将要与人共融时，应居住在最古老和最先的人，即\Person{亚当}内。因此，当祂更新那些从起初就有的事物，并由童贞女藉圣神重新塑造它们时，祂塑造了与起初同样的东西。那时大地仍是处女而未开垦，天主取泥土，从中形成理性的受造物，没有种子。\par

% structure: p0011 source=h2 target=subsection reason=relative-heading-rank; base=h2

\subsection{第五章 论\Person{耶肋米亚}的一段经文}

然而此处我可引用先知\Person{耶肋米亚}作为可靠而清晰的见证，他如此说：\BibleQuote{我下到陶匠的家里，看哪，他在轮盘上做了一件工。他用泥土做的器皿在陶匠手中坏了，他就按陶匠眼中看为好，又另做了一件器皿。} \BibleRef{耶 18:3} 因为\Person{亚当}由泥土形成时，仍是柔软潮湿的，尚未像瓦片一样变得坚硬不朽坏，罪就毁了他，如水一样流下滴在他身上。因此天主重新湿润他，重新塑造同一泥土归于祂的荣耀，先在童贞女的子宫中使它坚硬固定，并与圣言结合混合，将它带入生命，不再柔软破碎；免得它再次被外来的败坏之流冲溢而变软，并灭亡，正如主在祂的教导中，在寻找羊的比喻里所显示的；我主对旁边的人说：\BibleQuote{你们中间谁有一百只羊，失去一只，不把这九十九只撇在旷野，去找那失去的，直到找着呢？找着了，就欢欢喜喜地扛在肩上；回到家里，请朋友邻舍来，对他们说：你们和我一同欢喜吧，因为我失去的羊已经找着了。}\par

% structure: p0013 source=h2 target=subsection reason=relative-heading-rank; base=h2

\subsection{第六章 属灵羊群的总数；人是继天使之后的第二唱诗班，赞美天主；失羊比喻的解释}

既然祂是且确实是，在起初就与天主同在，且就是天主，\BibleRef{若 1:1}祂是天上万有的首领和牧者，所有理性的受造物都服从并服侍祂，祂有序地牧放和数点众多蒙福天使。因为这是不朽受造物的等同而完美的数目，按他们的种族和宗族分开，人也被收录在这群中。因为他也是无败坏而被造的，使他能尊荣万物的君王和创造者，回应从天上而来的悦耳天使的欢呼。但当发生因违反（天主的）诫命，他遭受了可怕而毁灭性的堕落，由此陷入死亡状态，为此主说祂从天上进入了（人类的）生命，离开了天使的等级和军队。因为山意指诸天，九十九只羊意指首领和掌权者，是元帅和牧者下去寻找迷失者时所撇下的。因为人必须被纳入这个目录和数目，主将他举起并包裹起来，使他不再，如我所说，被欺骗的波涛冲溢和吞没。因为圣言为此目的取了人的本性，为战胜那蛇，亲自摧毁那与人毁灭一同产生的定罪。因为理当如此：那恶者应由别者战胜，而应由那被他欺骗并夸口使其臣服者战胜，因为别无他法可毁灭罪和定罪，除非那同一人，就是曾被告说\BibleQuote{你是尘土，仍要归于尘土} \BibleRef{创 3:19} 的那人，应被重新创造，并撤销那因他而临到万有的判决，以致\BibleQuote{在亚当内}起初\BibleQuote{众人都死了，照样}再\BibleQuote{在基督内}，祂取了亚当的\Emph{本性和地位}，也要\BibleQuote{众人复活} \BibleRef{格前 15:22}。\par

% structure: p0015 source=h2 target=subsection reason=relative-heading-rank; base=h2

\subsection{第七章 基督的事工，属于天主也属于人，那独一者的事工}

现在我们似乎已对人如何成为独生子的器皿和外衣，以及那来居住在他内的祂是谁，说得足够多了。但关于没有\Emph{道德的}失衡或不和的事实，可以再次从开头简要思考。因为那位说以下话的人说得很好：那在其本性上是善的、正义的、圣的，其他事物因分有它而成为善；智慧与天主相连；另一方面，罪是不圣的、不义的、邪恶的。因为生命与死亡、朽坏与不朽坏，是彼此极度对立的两种事物。因为生命是\Emph{道德的}平等，而朽坏是不平等；正义与明智是和谐，而不义与愚昧是不和。现在，人处于这些之间，既非正义本身，也非不义；而是被置于不朽坏与朽坏之间，无论他倾向哪一方，都被说成分享那抓住他的东西的本性。当他倾向朽坏时，他就成为朽坏的和可死的；当他倾向不朽坏时，他就成为不朽的和不死的。因为，他被置于生命树与知善恶树之间，他尝了后者的果实\BibleRef{创 2:9}，就变成了后者的本性，他自己既非生命树，也非朽坏树；但他被显示为可死的，是由于他与朽坏的共融和同在；而又成为不朽的和不死的，是由于与生命的联系和共融；正如保禄也教导说：\Quote{朽坏不能承受不朽坏，死亡也不能承受生命}，他正确地界定朽坏和死亡是那朽坏和杀害的，而非那被朽坏和死亡的；不朽坏和生命是那赐予生命和不死的，而非那接受生命和不死的。因此，人既非不和与不平等，也非平等与和谐。但当他接受了不和，即过犯与罪，他就成为不和谐的和不雅的；但当他接受了和谐，即正义，他就成为和谐而雅致的器皿，为使那战胜死亡的主——不朽坏者——能使复活与肉体和谐，不让它再被朽坏所继承。关于这一点，这些陈述也已足够。\par

% structure: p0017 source=h2 target=subsection reason=relative-heading-rank; base=h2

\subsection{第八章 智慧的骨与肉；形成属灵的厄娃的肋骨，即圣神；女人——亚当的助手；与基督订婚的童贞女。}

因为已经藉着圣经中不容轻视的论证确立了：第一个人可以恰当地被指涉到基督本身，而不再是独生子的预象、代表和肖像，而是实际上成为智慧和圣言。\par

人，如同水一样，由智慧与生命所构成，已与倾注于他的那纯洁之光成为同一。因此宗徒直接将那论及亚当的话应用于基督。因为这样必能确定教会是由祂的骨与肉形成的；为此，圣言离开天上的父，降来\Quote{与祂的妻子结合}\BibleRef{弗 5:31}，并在祂受难的神魂超拔\OriginalTerm{神魂超拔}{trance}中睡眠，甘愿为她受死，为将教会呈现给自己，荣耀而无瑕，藉着洗濯\OriginalTerm{洗濯}{laver}使她洁净\BibleRef{弗 5:26}，以接受那属灵而蒙福的种子——这种子由祂以低语种在心灵的深处，并由教会如妇女一样孕育成形，以生育并滋养美德。同样，那\Quote{你们要生育繁殖}\BibleRef{创 1:18}的命令也得以圆满实现，教会因着圣言的结合与共融而日日增长，在伟大、美丽和众多上增进。圣言如今仍降临到我们中间，藉着对祂受难的纪念而陷入神魂超拔；否则，教会不能怀胎信徒，也不能藉重生的洗濯使他们重生，除非基督为了他们而空虚自己，好使他们能容纳祂——如我所说——藉着祂受难的重演，再次死去，从天降下，与教会——祂的妻子结合，并从自己肋旁取出一股力量，使所有在祂内被建立的人都成长，甚至那些藉洗濯重生的人，领受祂的骨和肉，即祂的圣德和光荣。因为谁若说智慧的骨与肉是理智与美德，便是最正确的；又说肋旁是真理之神、护慰者\OriginalTerm{护慰者}{Paraclete}，受光照者领受祂，便得以恰当地重生至不朽。因为任何人都不可能分享圣神，并被选为基督的肢体，除非圣言先降临到他身上并陷入神魂超拔，好使他充满圣神，并与那为他而躺卧睡眠者一同从睡眠中醒来，才能得到更新与恢复。因为按照先知的话\BibleRef{依 11:2}，祂可以恰当地被称为圣言的肋旁，即七重真理之神；天主在基督的神魂超拔中，也就是在祂的降生与受难之后，从祂那里取了一根肋旁，为祂预备了一个相称的助手\BibleRef{创 2:18}——我是指那些与祂订婚并嫁给祂的灵魂。因为圣经常常以教会的名称称呼信徒的聚会和群众，那些在进步中更完美的，被提升至成为教会的一个位格和身体。因为那些更优秀、更明白真理的人，脱离了肉身的邪恶，因着他们完美的净化与信德，成为基督的教会和相称的助手，如同童贞女许配并嫁给了祂，按照宗徒所说\BibleRef{格后 11:12}，这样他们领受祂教义的纯正真种子，就能与祂合作，在宣讲中帮助别人得救。而那些尚不完美、刚开始学习的，则被更完美的人如母亲般孕育至得救，并塑造他们，直到他们被生出、重生至美德的伟大与美丽；这样，他们本身在进步中成了教会，又协助生育和养育其他的儿女，在灵魂的容器中，如同在母腹中，完成圣言无瑕的旨意。\par

% structure: p0020 source=h2 target=subsection reason=relative-heading-rank; base=h2

\subsection{第九章 保禄宗徒的恩宠安排。}

现在让我们考虑著名的保禄的情况：当他在基督内还不完美时，他先被出生并哺乳，阿纳尼雅给他讲道，并在圣洗中更新他，正如宗徒大事录中的历史所述。但当他长成成人，并被建立起来，然后被塑造为属灵的完美，他便成了圣言的相称助手和新娘；他领受并孕育生命的种子，先前是孩童，如今成为教会和母亲，亲自为那些因他而相信主的人生产，直到基督也在他们内成形和诞生。因为他说：\Quote{我的小孩子们，我为你们再受生产之苦，直到基督在你们内成形}\BibleRef{迦 4:19}；又说：\Quote{我在基督耶稣内藉着福音生了你们}\BibleRef{格前 4:15}。\par

那么，显而易见，关于厄娃和亚当的陈述是指向教会和基督的。因为这确实是一个伟大的奥迹，而且是超自然的奥迹，我因软弱和愚钝，无法按其价值与伟大来谈论。然而，让我们尝试一下。我还要向你们讲述接下来的内容及其含义。\par

% structure: p0023 source=h2 target=subsection reason=relative-heading-rank; base=h2

\subsection{第十章 同一宗徒关于贞洁的道理。}

保禄在召唤众人归于圣化与纯洁时，将关于原祖和厄娃所说的话，以次要意义引用于基督与教会，为堵住无知者之口，使他们再无托辞。因为那些因情欲失控而不能自制的人，竟敢强解圣经，超出其真义，将\BibleQuote{你们要生育繁殖}\BibleRef{创 2:18}以及另一句\BibleQuote{为此人要离开自己的父亲和母亲} \BibleRef{创 2:24}曲解为放纵情欲的借口；他们不但不羞愧于违逆圣神，反而好像生来为此，挑旺潜伏的欲火，煽动激怒它。因此保禄断然割断这些不诚实的愚昧和编造的借口，在教导男人应如何对待妻子时，指明应如同基督对待教会一样，\BibleQuote{祂为自己舍弃了教会，为圣化她，以水的洗涤和圣言洁净她}。他回溯\BibleBook{}{创世纪}，提及关于原祖所说的话，并解释这些事与当前主题的关系，为杜绝那些借口生育而教导肉体官能满足的人对以上经文的滥用。\par

% structure: p0025 source=h2 target=subsection reason=relative-heading-rank; base=h2

\subsection{第十一章 同一论点。}

因为，童贞女们，请看他是如何竭尽全力地希望基督信徒保持贞洁，以多种论据向他们展示贞洁的尊贵，例如他说：\BibleQuote{关于你们信上所写的事：男人不接触女人是好的}，由此已非常清楚地表明不接触女人是好的，将此定规并明确地主张。但随后，他意识到不太能自制的人之软弱及其对结合的贪恋，便允许那些不能管束肉体的人使用自己的妻子，而不要因可耻地违反而陷入淫乱。在给予此许可后，他立刻补充了这些话\BibleRef{格前 7:5}：\BibleQuote{免得撒殚因你们不能自制而诱惑你们}；这意思是：\Quote{如果你们，像你们这样的人，因身体的不自禁和柔弱而不能完全自制，我宁肯允许你们与自己的妻子有往来，免得你们伪称完全自制而常受恶者诱惑，对他人妻子的贪欲炽燃。}\par

% structure: p0027 source=h2 target=subsection reason=relative-heading-rank; base=h2

\subsection{第十二章 保禄为寡妇及不与妻子同住者的榜样。}

来，现在让我们更仔细地审视眼前的这些话语，并注意宗徒并非无条件地将这些事赐予所有人，而是首先阐明导致他如此行事的理由。因为他先陈述了\BibleQuote{男人不接触女人是好的}，立刻补充说：\BibleQuote{但为了免于淫乱，男人当各有自己的妻子}\BibleRef{格前 7:2}——也就是说，由于你们不能约束自己的情欲而可能产生的淫乱——\BibleQuote{女人当各有自己的丈夫。丈夫要给予妻子应得的恩情，同样，妻子也要给予丈夫。妻子对自己的身体没有主权，丈夫却有；同样，丈夫对自己的身体没有主权，妻子却有。不要彼此亏负，除非两相情愿，暂时分房，为专务祈祷；以后再同房，免得撒殚因你们不能自制而诱惑你们。我说这话，是出于许可，并非出于命令。}\BibleRef{格前 7:2}这是经过深思熟虑的。他说\Quote{出于许可}，表明他是给予劝告，而非命令；因为他领受的是关于贞洁和不接触女人的\Emph{命令}，但对于那些不能，如我所说，节制自己欲望的人，则是\Emph{许可}。这些是关于与同一配偶结婚的男女，或将来要如此行的人而设立的；但现在我们必须仔细考察宗徒关于丧偶的男女，以及他就此所声明的话。\par

\BibleQuote{我说，}他继续说\BibleRef{格前 7:8}\BibleQuote{对未婚者和寡妇说：他们若能像我一样，守独处倒好。但如果不能自制，就让他们结婚，因为与其欲火中烧，不如结婚为好。}在这里他仍然坚持持守贞洁更好。因为他以自己为显著榜样，为激励他们效法，鼓励听者进入这种生活状态，教导一个已经从一夫一妻的束缚中释放的人，从此应当保持独身，正如他自己所做的那样。但如果另一方面，这对任何人来说有困难，因情欲冲动之强烈，他允许处于这样状况的人可以\Quote{出于许可}缔结第二次婚姻；并非因为他认为第二次婚姻本身是好的，而是判断它比欲火中烧更好。正如在预备复活节庆典的斋戒中，有人给一位病危的人食物，并说：说实话，我的朋友，你本宜勇敢地像我们一样坚持，并分尝同样的事物，因为今天甚至思考食物也是被禁止的；但既然你因疾病而软弱无力，无法承受，因此\Quote{出于许可}，我们建议你进食，以免因疾病而完全无力抵抗进食的欲望而死亡。宗徒在此也是这样说的：首先他说希望所有人都健康并节制，如同他自己一样；但随后允许那些受情欲疾病重压的人第二次结婚，以免他们完全被淫乱玷污，被性器官的搔痒驱策到混杂的交往中，认为这样的第二次婚姻远比焚烧和可耻之事更好。\par

% structure: p0030 source=h2 target=subsection reason=relative-heading-rank; base=h2

\subsection{第十三章 保禄关于童贞的教导的解释。}

关于节欲、婚姻、贞洁以及与人交往，以及其中哪些有助于在义德上进步，我现已说完应说的话；但关于童贞，若确有规定，仍需讨论。让我们也讨论这主题；因为经上这样记着：\BibleRef{格前 7:25} \BibleQuote{论到童身的人，我没有主的命令；但我蒙主怜悯成为忠信的，就提供我的意见。因着目前的艰难，我以为这样是好，就是说，人这样守身是好的。你有妻子缠着呢？不要寻求解脱。你没有妻子缠着呢？不要寻求妻子。你若娶妻，并不是犯罪；童女出嫁，也不是犯罪。然而这样的人肉身必受苦难；我却愿意你们免这苦难。}他极为谨慎地就童贞发表意见，并预备劝戒那愿意的人将他的童女出嫁，使那些有助于圣化的事都不是出于必要和强迫，而是出于心灵的自由意向。因为这样的祭献蒙天主悦纳，他不愿这些话作为权威和主的心意来说，尤其是关于将童女出嫁的事。因为在他说了\BibleRef{格前 7:28} \BibleQuote{童女出嫁，也不是犯罪}之后，紧接着极其谨慎地修正了言论，表明他是出于人的许可而非天主的许可来劝戒这些事。所以在他刚说了\BibleQuote{童女出嫁，也不是犯罪}之后，立即加上\BibleQuote{然而这样的人肉身必受苦难；我却愿意你们免这苦难。} \BibleRef{格前 7:28} 他的意思是：\Quote{我顾惜你们，像你们这样的，同意了这些事，因为你们愿意这样想，免得我显得是强逼你们，强迫任何人这样做。但若你们这些觉得贞洁难守的人，情愿转向婚姻，我认为约束自己在肉欲上是好的，不要使你们的婚姻成为滥用自己身体以致不洁的借口。}然后他接着说：\BibleQuote{弟兄们，我对你们说，时候减少了；从此以后，那有妻子的，要像没有妻子。}再后，继续勉励他们到同样的事情，他强调自己的话，大力支持童贞状态，并特别在前文之后加上下面的话，大声说：\BibleRef{格前 7:32} \BibleQuote{我愿你们无所挂虑。没有娶妻的，是为主的事挂虑，想怎样叫主喜悦。娶了妻的，是为世上的事挂虑，想怎样叫妻子喜悦。妇人和处女也有分别。没有出嫁的，是为主的事挂虑，要身体灵魂都圣洁；已经出嫁的，是为世上的事挂虑，想怎样叫丈夫喜悦。}现在，毫无疑问，谁都清楚：为主的事挂虑、讨天主喜悦，远胜于为世上的事挂虑、讨妻子喜悦。因为谁如此愚昧盲目，看不出这段话中保禄对贞洁的更高赞扬？\BibleQuote{我说这话，}他说，\BibleRef{格前 7:35} \BibleQuote{是为你们的益处，不是要牢笼你们，乃是要叫你们行合宜的事。}\par

% structure: p0032 source=h2 target=subsection reason=relative-heading-rank; base=h2

\subsection{第十四章　童贞是天主的恩赐：不可轻率立定守贞的意向。}

除此之外，还要思想他在上述话之后，如何称赞童贞是天主的恩赐。因此他拒绝了那些较为不能自制、却因虚荣而进到此状态的人，劝戒他们结婚，免得在他们精力旺盛之时，肉身激起欲望和情欲，驱使他们玷污灵魂。让我们细思他所立的规范：\Quote{若有人以为待他的童女不合宜，}他说，若她过了青春年华，事势需要，他可以随意而行，不算犯罪；让他们结婚吧。此处他恰当地将婚姻置于\Quote{不合宜}之上，指的是那些选择了童贞状态，后来觉得难以忍受、感到痛苦，口中在人前夸耀自己的坚忍，出于羞耻，但实际上已无力继续独身生活的人。但那些凭自己的自由意志和决心决定守身如玉、保持童贞纯洁的人，\BibleQuote{没有不得已的事} \BibleRef{格前 7:37} ，即情欲并未催逼他的腰肢去行房事——因为身体似乎各有不同——这样的人争战、奋斗，热切坚守自己的宣认，并出色地完成，他就劝勉这样的人坚持并保持，将最高的奖品授予童贞。因为他说，那能且渴望保持肉身纯洁的人，做得更好；那不能而依法进入婚姻、不暗中败坏的人，也做得好。关于这些事，现在说得够了。\par

谁愿意，可以拿起致格林多人的书信，逐一查考所有段落，然后思考我们所说的，相互比较，看是否完全和谐一致。这些，按我的能力，啊，\Person{Aretē}，作为我关于贞洁的献礼，献给你。\par

\Person{Euboulios}：历经许多事，啊，\Person{Gregorion}，她好不容易才进入正题，跋涉了浩瀚的语言海洋。\par

\Person{Gregorion}：似乎如此；但来吧，我必须按次序讲述其余的话，从头复述，趁我还似乎记得那声音在耳中回响，在它飞走消失之前；因为最近听到的事的记忆在年老者身上容易磨灭。\par

\Person{Euboulios}：那么请讲；因为我们得以享受这些讲论，真是愉快。\par

\Person{格雷戈里翁}：随后，正如你所观察到的，当\Person{塔莱娅}从她那平稳而连续的轨道下降到地面之后，\Person{德奥帕特拉}，她说，跟随着她，并说了以下的话。\par

\SourceRecord{Translated by William R. Clark.；Ante-Nicene Fathers；Vol. 6.；Edited by Alexander Roberts, James Donaldson, and A. Cleveland Coxe.；Buffalo, NY: Christian Literature Publishing Co.,；1886.；<http://www.newadvent.org/fathers/062303.htm>.}

% structure: p0001 source=h1 target=section reason=page-title

\section{\WorkTitle{十童女的婚宴（第四讲）}}

\Person{特奥帕特拉}\par

% structure: p0003 source=h2 target=subsection reason=relative-heading-rank; base=h2

\subsection{第一章 赞颂美德之必要——为那拥有能力者}

童女们，若演讲之术总沿同途、循同一径而行，则持续援引已陈之论者，必不能免于使你们厌倦。然而，论辩之流与途有万千之多，天主\Quote{多次并以多种方式}启迪我们，谁能选择退缩或畏惧呢？因为那领受恩赐者，若不以赞词装点可敬之事，便难逃罪责。来罢，我们也按各自的恩赐，歌颂基督最明亮、最荣耀的星辰——贞洁。因为这条圣神之路极其宽广。因此，让我们从那能就当前题目说出合宜恰当之言的起点开始，由此加以思考。\par

% structure: p0005 source=h2 target=subsection reason=relative-heading-rank; base=h2

\subsection{第二章 天主赐予人贞洁与童贞的保护，使人能脱离恶习的泥沼}

至少在我看来，似乎没有什么比贞洁更能使人复返乐园、转为不朽、与天主和好，并引导我们进入生命，成为人的救恩途径了。我现在要尽力说明我为何对此作如是观，好让你们清楚听见这已述恩宠的能力后，知道它赐予了我们何等伟大的福祉。古人，在人堕落、因过犯被逐之后，腐败的洪流大量倾泻，奔腾汹涌，不仅猛烈地冲走一切从外触及之物，而且涌入内部，淹没人的灵魂。人不断遭受此患，便哑然愚钝地被冲走，忽略掌舵，因无可把持之处。关于这些事，那些博学者曾言：灵魂的感官，当受到外来的情欲刺激侵袭，接住涌入的愚昧波涛的突袭而昏暗时，便使整个本性易导的船偏离天主之道。因此，天主怜悯我们处于如此境地、既不能站立也不能兴起，便从天降下最佳最荣耀的帮助——童贞，使我们借此把身体如船般系稳，得享平静，安然抵达锚地，正如圣神所证。因为在第一百三十六篇圣咏中，那些被基督握住并提起与祂同行于天上的灵魂，向天主欢唱感恩的诗歌，免得被世界和肉身的洪流淹没。为此，他们也说法郎在埃及是魔鬼的预像，因他残忍地命令将男婴丢入河中\BibleRef{出 1:16}，却保留女婴性命。因为那从亚当至摩西统治着这大埃及——世界的魔鬼\BibleRef{罗 5:14}，悉心将灵魂那属男性、理性的后代冲走并消灭于情欲的洪流，却渴望属血肉、非理性的后代增多繁殖。\par

% structure: p0007 source=h2 target=subsection reason=relative-heading-rank; base=h2

\subsection{第三章 解释大卫的那段经文；挂在柳树上的琴象征什么；柳树是贞洁的象征；溪水边的柳树}

但且不离开我们的主题，来罢，让我们手捧并审视这篇圣咏，那纯洁无瑕的灵魂向天主歌唱道：\Quote{我们曾在巴比伦的河边坐下，一追想熙雍就哭了。我们把琴挂在那中间的柳树上}，显然称琴为身体，他们将其挂在贞洁的树枝上，系于木头，免得被不洁的洪流夺去再冲走。因为巴比伦——解作\Quote{扰乱}或\Quote{混乱}——象征此生，四周有水环绕，我们坐在其中，水在我们四周流动，只要我们还在世上，邪恶的河流就不断冲击我们。因此，我们常怀恐惧，叹息哭泣向天主呼求，免得我们的琴被快感的波涛夺去，从贞洁之树上滑落。因为神圣著作处处以柳树为贞洁的象征，因它的花浸入水中，若饮下，便能熄灭那在我们内点燃感官欲望与情欲的一切，直至使其全然不结果实，使一切生育的倾向归于无效，正如荷马也为此原因称柳树为毁坏果实者。在依撒意亚书中，义人被称为\Quote{如溪水旁的柳树}\BibleRef{依 44:4}。诚然，童贞的嫩芽被举到荣耀的高处，是当义人——即那蒙赐保存并培育它者——以智慧浇灌它，被基督最温柔的溪流浸润之时。因为正如此树本性因水而发芽生长，童贞也因言语的滋润而开花成熟，使人能将身体挂在其上。\par

% structure: p0009 source=h2 target=subsection reason=relative-heading-rank; base=h2

\subsection{第四章 作者继续解释同一经文}

如果巴比伦的河流，正如贤者所言，是那令人心神混乱不安的肉欲之流，那么柳树就必是贞洁，我们可以将那些过分沉重、压向心灵的欲望器官悬吊并拉引其上，使它们不致被放纵的激流冲倒，像虫子一样被拖向不洁与腐朽。因为天主将童贞赐予我们，作为对不朽最有益、最有效的帮助，派遣它作为盟友援助那些为熙雍而奋斗、切望熙雍的人，正如圣咏所示——熙雍是灿烂的爱德及其诫命，因为\Quote{熙雍}被解释为\Quote{瞭望台的诫命}。现在，让我们在此列举后续的内容。因为为何灵魂们声称，那些俘掳它们的人要它们在外邦地歌唱上主的歌？确然，因为福音教导了一首神圣而隐密的歌，那是罪人与奸淫者向恶者唱的。因为他们侮辱诫命，履行恶神的旨意，把圣物扔给狗，把珍珠丢在猪前 \BibleRef{玛 7:6}，正如先知愤慨地论及那些人所说的：他们在外面读法律。原来犹太人不可走出耶路撒冷城门或自己的家门去读法律；因此先知严厉责备他们，宣告他们应受定罪，因为当他们违犯诫命、对天主行为不敬时，却装模作样地诵读法律，仿佛虔诚地遵守其规条；但他们并未在心中以信德牢牢接受它，而是拒绝它，以自己的行为否认它。因此他们在异乡唱上主的歌，以扭曲和贬低的方式解释法律，期待一个肉性的王国，寄望于圣言所说这个必将逝去的异世 \BibleRef{伯前 2:10}，在那里，俘掳他们的人以享乐诱惑他们，潜伏着要欺骗他们。\newline{}\par

% structure: p0011 source=h2 target=subsection reason=relative-heading-rank; base=h2

\subsection{贞女之恩，装饰她们以呈献给唯一新郎基督。}

如今，那些向愚昧之人歌唱福音的人，似乎是在异乡唱上主的歌，其土地并非基督所耕种；但那些穿戴并闪耀着最纯洁、光明、无掺杂、虔诚且合宜的童贞装饰，且对不安与痛苦的情欲毫无产出、不结果实的人，并不在异乡唱歌；因为他们的希望不导向那里，也不固着于可朽身体的欲望，亦不以低下的眼光看待诫命的意义，而是以高尚高贵的胸怀，仰视那至高的应许，渴慕天堂作为宜人的居所。因此天主赏识他们的心意，起誓应许赐予他们特选的荣耀，将他们安置并设立于\Quote{祂至高的喜乐之上}；因为祂如此说：\Quote{若我忘记你，耶路撒冷！愿我的右手丧失能力。我若不怀念你，不以耶路撒冷为最大喜乐，愿我的舌头贴住上膛。}此处以\Quote{耶路撒冷}意指，如我所言，这些无玷、不朽的灵魂，她们以克己之心、用无染的嘴唇饮下童贞的纯液，被\InnerQuote{许配给一位丈夫}，好能作为\InnerQuote{贞洁的童女呈献给基督} \BibleRef{格后 11:2}在天上，\InnerQuote{她们得了胜利，为无玷的奖赏而奋斗} \BibleRef{智 4:2}。因此，依撒意亚先知也宣告说：\Quote{起来，发光吧！因为你的光明已来到，上主的荣耀已照耀在你身上。} \BibleRef{依 60:1}这些应许，显然人人皆知，将在复活后实现。因为圣神并非指犹太地那个著名的城镇，而是指那天上的城邑，有福的耶路撒冷，祂宣告那是灵魂的集会——天主明确应许在崭新的安排中，将她们置于\Quote{祂至高的喜乐之上}，使那些身穿童贞最白袍衣的人，定居于不可接近之光的纯净居所；因为她们无意脱下自己的婚宴外衣——即，不以游思杂念松懈心志。\newline{}\par

% structure: p0013 source=h2 target=subsection reason=relative-heading-rank; base=h2

\subsection{处处时时当勤修并赞美童贞。}

再者，耶肋米亚所说的\Quote{处女不会忘记她的装饰，新娘不会忘记她的衣裳}，表明她不应因诡计与分心而放弃或松开贞洁的带子。因为心，正确而言，指的是我们的心思与意念。那束紧并稳固灵魂朝向贞洁之决心的胸带，便是对天主的爱——我们的元帅与牧者耶稣，也是我们的统治者与新郎，啊，卓越的贞女们，祂命令你们和我紧握它，使之不破裂、完好封存，直至终末；因为人不易找到比这一财富更大的帮助，它令天主喜悦和欣悦。因此，我说，我们人人都当操练并尊崇贞洁，时刻勤修并赞美它。\newline{}\par

\Person{阿瑞忒}啊，请接受这些我话语的初果，作为我教养与热忱的证明。\newline{}\Quote{我接受这礼物，}她转述\Person{阿瑞忒}回答说，\Quote{并请\Person{塔露莎}在你之后发言；因我需从你们每一位得到一篇发言。}她接着转述，\Person{塔露莎}稍作停顿，仿佛在心中思忖片刻，这样说道：\newline{}\par

\SourceRecord{Translated by William R. Clark.；Ante-Nicene Fathers；Vol. 6.；Edited by Alexander Roberts, James Donaldson, and A. Cleveland Coxe.；Buffalo, NY: Christian Literature Publishing Co.,；1886.；<http://www.newadvent.org/fathers/062304.htm>.}

% structure: p0001 source=h1 target=section reason=page-title

\section{《十贞女的筵席》（第五讲）}

\Person{塔路撒}\par

% structure: p0003 source=h2 target=subsection reason=relative-heading-rank; base=h2

\subsection{第一章　贞洁的奉献是伟大的恩赐}

我祈求你，\Person{亚肋特}，现在也赐予我帮助，使我所说的首先堪当配得上你，其次配得上在场诸人。因为我深信，从圣经中彻底学到，人所能献给天主的最大最光荣的祭献和恩赐，无可比拟的，就是贞女的生活。因为虽然在法律中，许多人按所许的愿完成了许多可赞叹的事，但唯独那些自愿奉献自己的人，才被称为完成了大愿。因为经文这样记载：\BibleQuote{主曾对梅瑟说：你告诉以色列子民说：无论男女，若愿特敬……奉献给主。} 一人许愿来到时献上金银器皿为圣所之用，另一人许愿献上果实的什一，另一人献上财产，另一人献上羊群中最好的，另一人奉献自己的生命；没有人能向主许一个大愿，除非他将自己完全奉献给天主。\par

% structure: p0005 source=h2 target=subsection reason=relative-heading-rank; base=h2

\subsection{第二章　亚巴郎献上三岁的母牛、母山羊和公绵羊的含义：每一年龄阶段都应奉献给天主；三重更与我们的年龄}

贞女们，我必须努力通过正确的解释，按经文的含义向你们阐明圣经的意思。那部分警醒自制、部分分心游荡的人，并非完全交付给天主。因此，完美的人必须献上一切，包括灵魂和肉身的事，好使他完整无缺。因此天主也命令\Person{亚巴郎}：\BibleQuote{你为我取一头三岁的母牛，一只三岁的母山羊，一只三岁的公绵羊，一只斑鸠和一只雏鸽。} 这话说得极妙；因为请留意，关于这些，他也命令说：把它们带到我这里，使它们免受轭的束缚——你的灵魂如母牛般无损伤，你的肉身如山羊般，因它行走于高峻险陡之处，以及你的理智如公绵羊般，好使他决不跳跃、坠落或滑离正道。因为当你将你的灵魂、你的官能和你的心思献给我时，你，\Person{亚巴郎}，就如此成为完美无缺的。他以三岁的母牛、母山羊和公绵羊象征这些，仿佛它们代表对天主圣三的纯洁认识。\par

也许他也象征我们生命和年龄的开始、中间和结束，愿意人们尽可能地纯洁地度过童年、成年和晚年，并将它们献给他。正如我们的主耶稣基督在福音中所命令的，这样指示：\BibleQuote{不要熄灭你们的光，也不要松开你们的腰带。因此你们也要像等候主从婚筵回来的人一样，当他来到敲门时，你们能立刻给他开门。你们是有福的，当他使你们坐下，并前来服侍你们时。如果他在第二更或第三更来到，你们也是有福的。} 因为请思量，贞女们，当他提到夜间的三更和他的三次来临，他象征性地预示我们人生的三个阶段：少年、壮年和老年；这样，如果他来临并在我们度过第一个时期——即我们仍是少年时——将我们带离世界，他就能接纳我们准备就绪、纯洁无瑕；第二和第三时期也是如此。因为晚更就是人萌芽和青年时期，那时理智开始因生活的变化而纷乱、蒙蔽，肉身增强力量驱使他纵欲。第二更是后来进入壮年时，他开始获得稳定，抵抗情欲和自负的骚动。第三更，则是多数想象和欲望消退，肉身萎缩而进入老年。\par

% structure: p0008 source=h2 target=subsection reason=relative-heading-rank; base=h2

\subsection{第三章　最宜从小修德}

因此，我们应当在心中点燃不可熄灭的信德之光，用纯洁束紧腰身，警醒并永远等候主，这样，如果他愿意来到并在人生第一时期、或第二、或第三时期将我们中的任何人带走，发现我们最为预备妥当，并履行他所委派的事，他就能使我们安眠在\Person{亚巴郎}、\Person{依撒格}和\Person{雅各伯}的怀中。\Person{耶肋米亚}说：\BibleQuote{人最好在青年时背负他的轭}， \BibleRef{哀 3:27} \BibleQuote{使他的灵魂不离开主。} 的确，从少年时就将颈项置于天主圣手之下，直至老年也不甩开那以纯洁心灵指引的骑手，实为善，因为恶者始终将心思拖向更恶的事。有谁不通过眼、耳、味、嗅和触，接受快乐与欢愉，以致不能忍受那作为驭手的节制之控制？节制如同驭手，勒紧并极力阻止马匹作恶。另一人若转念他事，便会各有见解；但我们说，那自幼努力保持肉身无玷、实践贞洁的人，是完美地将自己奉献给天主；因为贞洁迅速为追求它的人带来伟大而极向往的希望恩赐，干涸灵魂中败坏的贪欲和情念。但来吧，让我们解释我们如何将自己交付给主。\par

% structure: p0010 source=h2 target=subsection reason=relative-heading-rank; base=h2

\subsection{第四章　完全奉献与委身于天主：其意义}

户籍纪中\BibleQuote{发大愿}的记载，经稍加解释后便可证明，贞洁是超越一切愿的伟大誓愿。因为当我不仅努力保持肉身不受交合的玷污，也竭力使它不受其他不端行为的沾染时，我才真正完全奉献于主。因为经上说：\BibleQuote{未出嫁的女子，挂虑主的事，一心使主喜悦}；这不仅是为了在德行上不残缺而赢得部分光荣，而是按宗徒所言，在两方面都成全，使她肉身和灵魂都成圣，将她的肢体奉献给主。现在且说如何才算将自己完全奉献于主。例如，若我在某些事上开口，在其他事上闭口；若我开口是为按能力在真信德和应有尊敬下解释圣经、赞美天主，而闭口是给愚昧的言论安上门和锁，那么我的口便是纯洁的，奉献给天主。\BibleQuote{我的舌头是速记员的笔}，是智慧的器官；因为圣言藉着它，从圣经的深度和力量中写出最清晰的文字，甚至写下那亘古快写的记录者——主，祂急速而飞快地记录并实现圣父的旨意，听到\BibleQuote{快抢掠、急夺取}的话。对这样一位书记，可应用这句话：\Quote{我的舌头是速记员的笔}；因为一支优美的笔被祝圣并奉献给祂，写着比那证实人类学说的诗人和演说家更为优美的事物。此外，若我使自己的眼睛不贪恋肉体的魅力，不因不雅景象而喜悦，而是仰望天上的事，那么我的眼睛便是纯洁的，奉献给主。若我闭耳不听毁谤和诬蔑，只开启听天主圣言，与智慧人往来，那么我的耳朵便奉献给主。若我使自己的手远离不义之举、贪婪和放纵之行，那么我的手便对天主保持纯洁。若我约束脚步不走邪路，那么我的脚便奉献给主，不是走向公众场所和恶人聚集的宴席，而是走向正路，履行一些诫命。此外，若我也保持心灵纯洁，将所有心思奉献给天主；若我不思想邪恶，若愤怒和怒气不得胜于我，若我日夜默思上主的法律，那么我还剩下什么呢？这就是保持伟大的贞洁，发一个伟大的誓愿。\par

% structure: p0012 source=h2 target=subsection reason=relative-heading-rank; base=h2

\subsection{第五章 贞洁之愿及其法律中的礼仪；葡萄树、基督与魔鬼。}

如今，童贞女啊，我将尽力向你们解释其余的规定；因为这与你们的职责相关，包含关于贞洁的法律，教导我们应如何禁欲与如何进展至贞洁。因为经上如此记载：\BibleRef{户 6:1}\BibleQuote{上主训示梅瑟说：你告诉以色列子民说：无论男女，若许了一个特敬上主的愿，愿作一个纳齐尔人，他应戒饮清酒和烈酒，也不可饮清酒或烈酒酿的醋，不可喝任何葡萄汁，也不可吃鲜葡萄或干葡萄；在他奉献期内，凡葡萄树所产的，从葡萄粒到葡萄皮，都不可吃。}这意味着，凡献身并奉献于主的人，不可取用那罪恶植株的果实，因其本性容易导致迷醉和心神散乱。因为我们从圣经中感知到两种葡萄树，它们彼此分离，互不相同。一种结出不朽和义德，另一种结出疯狂和愚妄。那清醒而喜乐的葡萄树，从它的枝条上像教训般垂下一串串恩宠的果实，蒸馏出爱，便是我们的主耶稣，祂明确地对宗徒们说：\Quote{我是真葡萄树，你们是枝条；我父是园丁。}而那野性而致死的葡萄树，便是魔鬼，它滴下狂怒、毒汁和忿怒，正如梅瑟所记载的：\BibleRef{申 32:32}\BibleQuote{他们的葡萄树来自索多玛的葡萄树，来自哈摩辣的田地；他们的葡萄是毒葡萄，他们的葡萄串是苦的；他们的酒是龙的毒汁，是毒蛇的凶狠毒液。}索多玛的居民从这葡萄树上摘取葡萄，便被煽动起反天性的、无益的对男子的欲望。同样，在诺厄时代，人们因醉酒而陷入不信，被洪水淹没而溺死。加音也从这树取汁，玷污了他杀兄弟的手，以亲族的血污秽了土地。同样，外邦人因醉酒而激发起杀伐之战的激情；因为人喝酒所受的刺激和迷途，远不如因愤怒和怒气而严重；人因酒而醉酒迷途，也不如同因忧伤、爱情或不贞而严重。因此，命令童贞女不可品尝这葡萄树的果实，为使她清醒谨慎，脱离生活的挂虑，并为圣言点燃义德之光明的火炬。主说：\BibleRef{路 21:34}\BibleQuote{你们要谨慎，免得你们的心因沉湎于宴饮、醉酒和生活的挂虑而沉重，那日子如同罗网般突然临到你们身上。}\par

% structure: p0014 source=h2 target=subsection reason=relative-heading-rank; base=h2

\subsection{第六章 烈酒——一种人工仿制且仍致醉的酒；童贞女应远离与罪相近之事；香坛（童贞女的象征）。}

此外，不仅禁止贞女以任何方式触摸由那葡萄所出之物，就连那些与它们相似且相近之物也禁止。因为人工酿制的\OriginalTerm{希克拉}{Sikera}，无论由棕榈还是其他果树制成，都被称为一种仿制的酒。因为正如酒饮能倾覆人的理性，这些也极其如此；坦白地说，智者惯于将一切除酒之外导致醉酒和心神涣散的东西称为希克拉之名。因此，为使贞女在防范那些本性为恶的罪时，不被那些与它们相似且相近之物玷污，即克服一种却被另一种克服——亦即以不同布料的编织品、或宝石与黄金及其他身体装饰来装饰自己，这些事物使灵魂迷醉——故有此命令，要她不可沉溺于女性的软弱与嬉笑，激起自己趋于诡诈与愚妄言谈，这些使心思旋转混乱；正如在另一处所指示的：\Quote{不可吃鬣狗及与它相似的动物；也不可吃鼬鼠及这类受造物。}\par

此外，传承下来的教导是：天主的无血祭台象征着贞洁者的集会；如此，贞洁似乎是一件伟大而光荣的事。因此，它应当被保存得毫无玷污、全然纯洁，不与肉体的不洁有任何关联；它应被安置在约证之前，用智慧镀金，为至圣所，向主发出爱德的馨香；因为祂说：\BibleRef{出 30:1} \BibleQuote{应用皂荚木做一座烧香的坛。要用皂荚木做杠，用金包裹。要把坛放在法柜前的幔子外，对着约证上的施恩座，就是我要与你相会的地方。亚郎要在坛上烧芬芳的香，每早晨整理灯的时候，要烧这香。黄昏点灯的时候，亚郎也要烧这香，作为世世代代在主面前常烧的香。在这坛上不可献异样的香，不可献燔祭、素祭，也不可奠上奠酒。}\par

% structure: p0017 source=h2 target=subsection reason=relative-heading-rank; base=h2

\subsection{第七章 在法律的阴影与天堂的实在之间的教会}

如果法律按宗徒所说，是属神的，包含\Quote{未来美物的影像}，那么来吧，让我们揭去覆盖其上的文字幔子，思考它赤裸而真实的含义。希伯来人奉命装饰会幕作为教会的预象，使他们能通过可感知的事物，预先宣告神圣事物的影像。因为西奈山上向梅瑟所显示的样式\BibleRef{出 25:40}，他铸造会幕时应依照它的，乃是天上居所的精确写照，我们现在通过预象比从前更清晰地感知它，却比我们亲眼看到实在更加模糊。因为就我们目前的状况而言，纯粹的真理尚未与世人混合，世人在此无法承受纯粹不朽的视线，正如我们无法直视太阳的光芒。犹太人宣称那（借给他们）的天上之物的影像的阴影，是与实在相隔第三层；但我们则清晰可观天上秩序的影像；因为真理将在复活后准确彰显，那时我们将看到天上的会幕（天上的城\BibleQuote{其设计者和建造者乃是天主} \BibleRef{希 11:10}）\BibleQuote{面对面}，而非\BibleQuote{模糊}和\BibleQuote{部分}\BibleRef{格前 13:12}。\par

% structure: p0019 source=h2 target=subsection reason=relative-heading-rank; base=h2

\subsection{第八章 双重祭台，寡妇与贞女；金子是贞洁的象征}

如今，犹太人预言了我们的状况，而我们则预言天上的状况；因为会幕是教会的象征，教会又是天堂的象征。因此，既然这些事情如此，并且会幕被视为教会的预象，如我所说，那么理所当然祭台应象征教会中的某些事物。我们已将铜祭台比作寡妇的团体和巡回；因为她们是天主的活祭台，她们将牛犊、什一之物和自愿祭献作为牺牲献给主；但至圣所内、约证之前的金祭台，其上禁止献祭和奠酒，则指向处于贞洁状态的人，因为她们的身体如纯金般保存纯洁，远离肉体的结合。金子因两个理由受称赞：第一，它不生锈；第二，它的颜色在某种程度上似乎类似太阳的光线；因此它适当地是贞洁的象征，不容许任何污点或斑点，而是常以圣言的光辉照耀。因此，它也\Emph{更接近天主}，在至圣所内，在幔子之前，以无玷之手，如同馨香，向主献上祈祷，蒙悦纳为芬芳之气；正如若望所指明的，说二十四位长老香炉里的香，乃是圣人们的祈祷。阿勒忒啊，这就是我即兴所能献给你的，关于贞洁的论述。\par

当塔罗萨说完这些话后，塞奥帕特拉说阿勒忒用她的权杖触了阿加特，而阿加特察觉后，立刻起身回答。\par

\SourceRecord{Translated by William R. Clark.；Ante-Nicene Fathers；Vol. 6.；Edited by Alexander Roberts, James Donaldson, and A. Cleveland Coxe.；Buffalo, NY: Christian Literature Publishing Co.,；1886.；<http://www.newadvent.org/fathers/062305.htm>.}

% structure: p0001 source=h1 target=section reason=page-title

\section{十贞女的婚筵（第六论）}

\Person{阿加特}\par

% structure: p0003 source=h2 target=subsection reason=relative-heading-rank; base=h2

\subsection{第一章 童贞永恒荣耀的卓越；灵魂按照天主子——天主肖像的肖像受造；魔鬼追求灵魂}

\Person{阿勒忒}，因着你能说服人并继续这一卓越谈论的巨大信心，若你与我同行，我也将竭尽所能，为当前议题的讨论贡献一些东西，一些与我自己力量相称、而无法与已言者相比的东西。因为我在哲学探讨中无法提出任何能与先前那些多样而出色完成的事物相竞争的东西。若我试图与智慧上高于我者匹敌，我将背负愚钝的责备。然而，若你仍愿容忍那些按自己能力说话的人，我将努力发言，至少不缺乏善意。现在让我开始。\par

诸位贞女，我们都带着一份与\Emph{神圣}智慧相关且亲密的非凡美丽来到此世。因为人的灵魂，当它反映出其肖像无瑕的代表，以及天主观看并塑造他们成为不朽不坏之形状的那面容的特征时，它们便最准确地相似那生了并形成了它们的那一位。因为无生无体的美丽，既不开始也不腐朽，而是不变、不老、一无所缺，祂安息于自身，在不可言说、不可接近之处的光中，以祂能力之圆环包容万有，创造并安排，按祂肖像的肖像造了灵魂。因此，灵魂也是合理的、不死的。因为如我所说，它被造为独生子的肖像，具有无比的美丽，因此恶灵爱它，并图谋和竭力玷污这神圣可爱的肖像，正如\Person{耶肋米亚}先知所示，责斥\Place{耶路撒冷}说：\BibleQuote{你额上有娼妓的相，你拒绝羞愧。} \BibleRef{耶 3:3}这是指那向侵犯她以玷污她的权力卖淫的她。因为她的情人是魔鬼及其天使，他们计划通过与灵魂交合来玷污和污染我们合理且明晰的美丽心灵，并渴望与每一个许配给主的灵魂同居。\par

% structure: p0006 source=h2 target=subsection reason=relative-heading-rank; base=h2

\subsection{第二章 十贞女的比喻}

因此，若有人保持这美丽不受侵犯、不受伤害，如同那建造者所形成、所塑造的一样，效仿那人所代表和相似之永恒可理解的性体，成为荣耀神圣的肖像，他将由此被转移到天堂，即蒙福者的城邑，并在那里如同在圣所中居住。当我们的美丽受到童贞的保护，不被外在腐败的热气玷污，而是保持在自身中，以义德装饰，作为新娘带给天主子时，它便最完美地保存为无瑕、完全的；正如祂自己提示，劝勉贞洁之光应在他们肉体内被点燃，如同灯一般；因为十贞女的数目象征那信仰\Person{耶稣基督}的灵魂，以十表示唯一通往天堂的正确道路。现在五人是谨慎而聪明的；五人则是愚拙不智的，因为他们没有预见用油装满自己的器皿，缺乏义德。通过这些，祂指那些努力达到童贞边界、竭尽全力实现这爱、行事有德操节制、并宣称以此为目标的灵魂；但他们轻视它，被世界的变迁所征服，反而成了美德模糊肖像的草图，而非代表活生生真理的工人。\par

% structure: p0008 source=h2 target=subsection reason=relative-heading-rank; base=h2

\subsection{第三章 追求童贞的同样努力与不同结果}

如今，当说到\EditorialNote{玛窦福音第二十五章}所记的\Quote{天国好比十个童女，拿着灯，出去迎接新郎}时，这意味着她们走上了同样的目标之路，正如标记X所表明的。她们因共同的志向而认同同样的终向，因此被称为十位，正如我已经说过的，她们选择了同样的志向；但尽管如此，她们出去迎接新郎的方式却不尽相同。因为有些人为她们用油点亮的灯准备了丰富的未来滋养，而另一些人则粗心大意，只顾眼前。因此，她们被分成相等的两个五，因为一类人保持了五种感官——大多数人认为那是智慧之门——纯洁无瑕，不受罪恶玷污；而另一类人则相反，被无数罪恶所败坏，以邪恶玷污了自己。因为她们约束了这些感官，使它们远离正义，从而结出更丰盛的过犯果实，结果是她们被禁止进入神圣的庭院，被拒之门外。因为无论我们是通过这些感官行善还是作恶，我们的善习或恶习都因此得以巩固。正如\Person{塔卢萨}所说，眼睛、耳朵、舌头等感官都有贞洁；同样，那保持五种德性通道——视觉、味觉、嗅觉、触觉和听觉——的信德完好无损的人，被称为五个童女，因为她将五种感官的形式纯全地保持给基督，如同灯一样，使圣洁之光从每一种中清晰地照射出来。因为肉体实际上就是我们那盏五灯的灯，灵魂在复活之日站在新郎基督面前时，将如同火炬般举着它，通过所有的感官展示出清晰而明亮的信德，正如祂亲自教导说：\Quote{我来是把火投到地上，我是多么希望它已经燃烧起来！}这火意味着我们的身体，祂希望祂教义的火速而炽热的运作在其中被点燃。如今，这油代表智慧和正义；因为当灵魂毫无保留地将这些倾注于身体时，德行之灯就不可熄灭地点燃了，使它的善行在人前照耀，从而光荣我们在天之父。\BibleRef{玛 5:16}\par

% structure: p0010 source=h2 target=subsection reason=relative-heading-rank; base=h2

\subsection{第四章 灯中的油意味着什么}

他们曾在肋未记中提供这类油，\BibleRef{肋 24:2}，\Quote{纯橄榄油，打成的，用来点灯，使灯在帷幔前……常燃不熄，在上主面前。}但以色列人奉命从晚上到早晨保持微弱的光。因为他们的光似乎类似先知的话，鼓励节制，并藉百姓的行为和信德得滋养。然而，圣殿（灯在其中常燃）指的是\Quote{他们继承的产业}，因为一盏灯只能在一所房屋中照耀。因此，必须在天亮前点燃。因为经上说，\BibleRef{肋 24:3}，\Quote{\Emph{他们应燃点它}直到早晨，}也就是说，直到基督的来临。但贞洁和正义的太阳已经升起，就不再需要\Emph{其他}光了。\par

因此，只要这百姓为灯储备滋养，以他们的行为供给油，节制的光就在他们中没有熄灭，而是一直照耀，在\Quote{他们继承的产业}中发光。但当油因他们从信德转向不节制而耗尽时，光就完全熄灭了，以至于童女们不得不再次借由彼此传递的光点燃她们的灯，从高处将不朽之光带到世上。那么，让我们现在就充足地供给善行和明智的油，清除一切会压垮我们的朽坏；免得在新郎迟延的时候，我们的灯也同样熄灭。因为这迟延是基督显现之前的间隔。如今，童女们的打盹和睡觉表示离世；半夜是假基督的王国，在此期间毁灭天使越过房屋。但当有人说\Quote{看，新郎来了，你们出去迎接他}时所发出的呼喊，是从天上听到的声音和号角，那时圣徒们，所有身体复活后，将被提升，驾云迎接主。\BibleRef{得前 4:16}\par

因为应当注意，\Emph{天主的}话说到，在呼喊之后，所有的童女都起来了，也就是说，死人将在那从天上来的声音之后复活，正如保禄也暗示的，\BibleRef{得前 4:16}，\Quote{主必亲自从天降下，有呼喊声，有总领天神的声音，也有天主的号角；那在基督内死了的人要先复活。}那就是帐幕（指身体），因为它们死了，被灵魂脱去。\Quote{然后我们这些活着的人，将和他们一同被提升，}意指我们的灵魂。因为我们真正活着的是灵魂，它们与身体（再次穿上）一起，将驾云迎接祂，提着修剪过的灯，不是带着任何外来的属世之物，而是像星星一样，放射着明智和节制的光，充满以太的光辉。\par

% structure: p0014 source=h2 target=subsection reason=relative-heading-rank; base=h2

\subsection{第五章 贞洁的奖赏}

啊，佳美的贞女们，这就是我们奥秘的庆典；这就是那些在贞洁中受训者所行的神秘仪式；这就是贞洁战斗的\Quote{无玷的报酬}\BibleRef{智 4:2}。我许配给了圣言，并从圣父那里获得不朽的永生荣冠和财富；我戴着智慧光明而不凋的花朵，在永恒中凯旋。我与基督一同在歌咏班中，祂在天上分施赏报，围绕着无始无终的君王。我成了不可接近之光的持炬者，并加入他们的行列，在总领天使的新歌中，彰显教会的新恩宠；因为圣言说，贞女的行列总是跟随主，并在祂所在之处与祂共融。这就是若望在提到十四万四千人时所象征的。\par

那么，去吧，你们这新时代的贞女队伍。去吧，用正义充满你们的器皿，因为时辰将至，你们必须起来迎接新郎。去吧，轻轻抛开那些迷惑、蛊惑灵魂的世俗诱惑与享乐；这样你们就能获得所应许的，\Quote{我凭那指示我生命之道的那位起誓。}这由先知们编织的冠冕，我从先知的花园中采来，献给你，啊，\Person{阿瑞忒}。\par

\Person{阿加忒}如此出色地结束了她的讲论，她说完后，得到了赞扬。\Person{阿瑞忒}又命\Person{普罗基拉}发言。她起身，走到入口前，这样说道。\par

\SourceRecord{Translated by William R. Clark.；Ante-Nicene Fathers；Vol. 6.；Edited by Alexander Roberts, James Donaldson, and A. Cleveland Coxe.；Buffalo, NY: Christian Literature Publishing Co.,；1886.；<http://www.newadvent.org/fathers/062306.htm>.}

% structure: p0001 source=h1 target=section reason=page-title

\section{十童女的宴会（第七讲）}

\Person{Procilla}\par

% structure: p0003 source=h2 target=subsection reason=relative-heading-rank; base=h2

\subsection{第一章. 什么是真实合宜的赞美方式；圣父大于圣子，不是在本体上，而是在次序上；贞洁是百合花；信实的灵魂与贞女，是唯一基督的唯一新娘。}

我不可以拖延，喔，\Person{Arete}，在如此讲论之后，因为我深信天主那丰富的智慧，祂慷慨而广泛地赐给祂所愿意的人。因为航海有经验的人说，同一阵风吹向所有航行的人；而不同的人，以不同方式掌舵，力图到达不同的港口。有些人顺风，有些人侧风；然而两者都轻易完成航程。同样地，那\Quote{聪敏的圣神，唯一神圣的}，从天上圣父的宝库中温柔地吹下，赐给我们清晰顺遂的知识之风，足以引导我们的话语无误。现在是我发言的时候了。喔，贞女们，这是唯一真实合宜的赞美方式：当赞美者提出一个比所有被赞美者更好的见证。因为由此可以确切知道，赞美并非出于偏袒、勉强或名声，而是依照真理和不谄媚的判断。因此，先知和宗徒们更详尽地谈论天主圣子，并赋予祂超越他人的神性时，并未将他们对祂的赞美归之于天使的教训，而是归于那一切权能和权力所依赖的那位。因为那在圣父以后大于万有的那位，理当以唯独大于祂自己的圣父为祂的见证\BibleRef{若 14:28}。所以，我不凭人的传闻来赞扬童贞，而是凭那眷顾我们、并承担一切的那位，祂表明自己是这恩宠的农夫、热爱其美者、以及合宜的见证。这在\WorkTitle{雅歌}中，任何愿意看的人都十分清楚，在那里基督亲自赞扬那些坚定持守童贞的人，说：\BibleRef{歌 2:2} \BibleQuote{我的爱卿在少女中，好像荆棘中的百合花。} 祂将贞洁的恩宠比作百合花，因它的纯洁、芬芳、甜美和喜乐。因为贞洁如同春天的花朵，总是从它白色的花瓣中轻柔地散发不朽。因此，祂不羞于承认祂爱它年青时的美丽，用以下话语：\BibleRef{歌 4:9} \BibleQuote{我的妹妹，我的新娘，你夺去了我的心；你以你的一只眼睛，你颈上的一条项链夺去了我的心。你的爱情多么美丽，我的妹妹，我的新娘！你的爱情比美酒更甘甜，你香液的芬芳超越一切香料！我的新娘，你的嘴唇滴流蜜汁，你的舌下有蜜有奶；你衣服的芬芳好似黎巴嫩的馨香。我的妹妹，我的新娘，是关闭的花园，是一座关锁的花园，是一个封锁的泉源。}\par

这是基督对那些已达到童贞边界的人所宣告的赞美，将它们全部归于祂新娘这一个名号下；因为新娘必须许配给新郎，并被称为祂的名。而且，她必须无玷无污，如同一座封闭的花园，其中生长着天上芬芳的一切香气，好让基督独自前来采摘，以无形体的种子盛开。因为圣言不爱任何肉体之事，因为祂的本性不会满足于任何可朽之物，如手、脸或脚；祂注视并喜悦那非物质和属灵的美，而不触碰身体的美。\par

% structure: p0006 source=h2 target=subsection reason=relative-heading-rank; base=h2

\subsection{第二章. 对\WorkTitle{雅歌}那段经文的诠释}

现在，喔，贞女们，请注意，当祂对新娘说：\BibleQuote{我的妹妹，我的新娘，你夺去了我的心，} 祂显示了悟性之明澈的眼目，当内在的人洗洁它，并更清楚地看见真理时。因为人人皆知，视力有双重力量：一为灵魂的，一为身体的。但圣言并不表示爱那肉体的，而只爱那悟性的，说：\BibleQuote{你以你的一只眼睛，你颈上的一条项链夺去了我的心；} 意思是：以你最可爱的思想之眼，你促使我的内心爱上，从内在放射出贞洁的光辉之美。颈上的项链是由各种宝石制成的；那些顾惜身体的人，将这可见的装饰戴在肉体的外在颈项上，以欺骗观看者；但那些贞洁生活的人，相反地，以真正由各种宝石制成的——即自由、宽宏、智慧和爱——装饰自己内在，很少关心那些暂时的装饰，它们如同短暂开花的叶子，随身体的变化而枯萎。因为人可见双重之美，主接纳那内在且不朽之美，说：\BibleQuote{你以你颈上的一条项链夺去了我的心；} 意思是要表明祂是因内在之人的光辉荣耀而被迫去爱，正如圣咏作者也见证说：\BibleQuote{君王的女儿全然荣耀内在。}\par

% structure: p0008 source=h2 target=subsection reason=relative-heading-rank; base=h2

\subsection{第三章. 贞女是基督同伴中首先的殉道者}

不要让任何人以为所有其余的信众都被判有罪，以为只有我们贞女才能被引领获得应许，而不明白将有支派、家族和等级，依据各人信德的类比。保禄也如此阐明，说：\BibleRef{格前 15:41}\Quote{太阳有一种光荣，月亮有另一种光荣，星辰也有另一种光荣，而且星辰与星辰的光荣也有分别。死人的复活也是这样。}主并不许诺将同样的荣誉赐予众人；而是向一些人应许他们将被列在天国里，向另一些人应许承受大地，向另一些人应许看见圣父。\BibleRef{玛 5:3}在这里，祂也宣布：贞女的队列和圣洁歌咏团将首先与祂一同进入新安排中的安息，如同进入新房。因为她们是殉道者，并非因一时忍受身体的痛苦，而是终生忍受，毫不畏缩地真正为贞洁的奖赏进行奥林匹克竞赛；她们抵抗着快乐、恐惧、忧愁及人类不义的其他邪恶的猛烈折磨，首先赢得奖赏，置身于领受应许者的更高行列。无疑地，这些灵魂就是圣言称为祂独选的配偶和祂的妹妹的，而其余的称为妾、贞女和女儿，如此说：\BibleRef{歌 6:8}\Quote{有六十王后，八十妾，并有无数的贞女。我的鸽子，我的完人，只有一个；她是她母亲的独一者，是生她者的宠儿：众女儿看见了她，就祝福她；王后和妾也看见了她，就赞美她。}因为教会的女儿显然众多，但在她眼中唯有一个是独选且最为珍贵的，就是贞女的等级。\par

% structure: p0010 source=h2 target=subsection reason=relative-heading-rank; base=h2

\subsection{第四章。经文解释；洪水前的王后，即圣善的灵魂；洪水后的妾，即先知们的灵魂；先知书中的神圣种子以产生属灵的后裔；圣言在先知中的婚配如同隐秘的。}

现在若有人对此有疑问，因这些事并未在任何地方完全阐明，并仍希望更充分地领会其属灵意义，即王后、妾和贞女是什么，我们将说，这些可能是论及那些从起初以来在时间的进程中因正义而杰出的人；如洪水前的和洪水后的，以及基督之后的。教会就是配偶。王后是洪水前的那些帝王的灵魂，他们中悦了天主，即围绕\Person{亚伯尔}、\Person{舍特}和\Person{哈诺客}的人。妾是洪水后的那些，即先知们的灵魂，在教会与主订婚之前，主与她们结合如同妾一般，将真道播种在不朽而纯洁的哲学中，使她们孕育信德，为祂产生救恩的圣神。因为基督与之交往的灵魂产生这样的果实，这些果实带有永世难忘的名声。因为如果你们查看\Person{梅瑟}、\Person{达味}、\Person{撒罗满}、\Person{依撒意亚}以及随后的先知们的书卷，啊，贞女们，你们会看到他们与天主之子交往留下了什么后代，以拯救生命。因此圣言以深刻的理解称先知们的灵魂为妾，因为祂没有公开聘娶她们，如同祂聘娶教会一样，并为了她宰杀了肥牛犊。\BibleRef{路 15:23}\par

% structure: p0012 source=h2 target=subsection reason=relative-heading-rank; base=h2

\subsection{第五章。六十位王后：为何是六十，为何是王后；第一个时代圣人的卓越。}

此外还有一件事需要考虑，以免我们遗漏任何必要之事：为什么祂说王后是六十位，妾是八十位，贞女多到不可胜数，而配偶却只有一个。首先让我们谈谈六十位。我想，祂以六十王后指代那些从第一个人开始一直到\Person{诺厄}为止相继取悦天主的人，因为这些人不需要规矩和法律来获得救恩，那时六日创世还近在眼前。因为他们记得天主在六日内创造了万物，以及乐园中所造之物；人如何因领受不可触摸知善恶树的命令而搁浅，邪恶之首诱惑了他。\BibleRef{创 3:3}由此，祂将象征性的名称六十王后赐给那些灵魂，他们从创世以来相继选择天主作为他们的爱慕对象，并且几乎可以说是第一个时代的后裔，是伟大六日工程的邻居，因为他们就像我说的，生于六日之后不久。这些人享有极大的尊荣，与天使交往，常常看到天主有形地显现，而非在梦中。你们想想，\Person{舍特}、\Person{亚伯尔}、\Person{厄诺士}、\Person{哈诺客}、\Person{默突舍拉}和\Person{诺厄}对天主有何等信心，他们是正义的首批热爱者，是写在天上的首生儿女中的首批，\BibleRef{希 11:23}被认为堪当天国，如同救恩植物的初果，作为早熟的果子献给天主。关于这些，说这些就够了。\par

% structure: p0014 source=h2 target=subsection reason=relative-heading-rank; base=h2

\subsection{第六章。八十位妾，是什么；降生成人的知识传达给先知们。}

关于那些妾室，仍需言说。对于洪水之后生活的人们，对天主的认识从此更加遥远，他们需要另一种教导来抵御邪恶，并成为他们的帮助者，因为偶像崇拜已经悄然潜入。因此，天主为使人类不致因忘却善事而完全毁灭，命令自己的圣子向先知们启示祂将来在世上藉肉身的显现；在此显现中，属灵的第八日之喜乐与知识将被宣告，它将带来罪的赦免与复活，并由此割除人的情欲与败坏。因此，祂将亚巴郎以来的先知行列称为八十个贞女，是因割损的尊贵，即包含数字八，法律也依此制定；因为他们，在教会与圣言订婚之前，首先接受了神圣的种子，并预言了属灵第八日的割损。\par

% structure: p0016 source=h2 target=subsection reason=relative-heading-rank; base=h2

\subsection{第七章 贞女、古义人；教会，唯一的新娘，超越其他}

现在，祂将那些属于不可胜数的会众、次于更优秀者、却实践了正义并以青春高尚的精力与罪争战的人称为贞女。但这些之中，无论是王后、妾室，还是贞女，都不能与教会相比。因为教会被视为完美而蒙选者，超越这一切，由所有宗徒组成并构造，是那位在青春和贞洁之美上超越一切的新娘。因此，她也蒙众人祝福和赞美，因为她自由地看见和听见了那些人甚至渴望看见片刻却未看见、渴望听见却未听见的事。因为我们的主对祂的门徒说：\BibleQuote{你们的眼睛是有福的，因为看得见；你们的耳朵是有福的，因为听得见。我实在告诉你们，许多先知曾渴望看见你们所看见的，却没有看见；渴望听见你们所听见的，却没有听见。}\BibleRef{玛 13:16}因此，先知们称他们有福，并钦佩他们，因为教会堪当分享那些他们未能听见或看见的事。因为\BibleQuote{有六十王后，八十妾室，和无数贞女。我的鸽子，我的完全人，只有一位。}\BibleRef{歌 6:8}\par

% structure: p0018 source=h2 target=subsection reason=relative-heading-rank; base=h2

\subsection{第八章 基督的人性，祂唯一的鸽子}

现在，谁能说不？新娘就是主无玷的肉身；为这肉身的缘故，祂离开了圣父，降临到这里，与它结合，并降生成人居住在其中。因此祂譬喻称它为鸽子，因为这生物温顺驯良，易于适应人的生活方式。因为可以说，唯有她被发现是无玷无瑕的，在正义的荣美上超越一切，以至于凡最完美地悦乐天主的人，在德性的比较中无人能近于她。因此，她堪当成为独生者王国的一份子，与祂订婚并联合。在第四十四篇圣咏中，那位从众多中被选出的王后，站在天主的右边，身披德性的金饰，君王的欲望了她的美貌，正如我所说，就是那无玷而有福的肉身，圣言亲自将她带入诸天，呈献在天主的右边，\BibleQuote{绣有各种色彩}，即追求不朽，祂象征性地称之为金穗。因为这件外衣是多彩的，由各种美德编织而成，如贞洁、明智、信德、爱德、忍耐及其他善行，它们遮盖了肉身的丑陋，以金饰妆饰人。\par

% structure: p0020 source=h2 target=subsection reason=relative-heading-rank; base=h2

\subsection{第九章 紧接王后和新娘之后的贞女}

此外，我们还必须进一步思考圣神在圣咏其余部分中传达给我们的内容，在圣言所取的人性被高举于圣父右边之后。祂说：\BibleQuote{那些跟从她的贞女，要陪伴她，并且要被带到你面前。她们要欢欣快乐地被带来，进入君王的宫殿。}现在，圣神似乎相当清楚地赞美了贞洁，次于主的新娘，正如我们已经解释过的；那位新娘应许贞女们将带着喜乐和快乐，在天使的守护和护送下，次于全能者前来。因为贞洁的荣耀确实是如此可爱而可羡，以至于仅次于那位被主高举、以无罪荣耀呈献于圣父面前的王后，贞女的歌队和行列陪伴着她，被分配给次于新娘的位置。阿勒忒啊，愿我这些关于贞洁的言论被刻在纪念碑上。\par

\Person{普罗奇拉}说完这番话后，\Person{特克拉}说：轮到我在她之后继续竞赛了；我欢喜，因为我也拥有恩赐的言词智慧，我感觉自己如同竖琴，内在已被调音，准备好以优雅且得体的话语发言。\par

\Person{阿勒忒}说：\Person{特克拉}啊，我最乐意欢迎你的准备，我信赖你，将依你的能力赐予我合适的讲论；因为你在普遍的哲学与教导上不亚于任何人，并受教于保禄，知道该怎样宣讲福音和神圣的教义。\par

\SourceRecord{Translated by William R. Clark.；Ante-Nicene Fathers；Vol. 6.；Edited by Alexander Roberts, James Donaldson, and A. Cleveland Coxe.；Buffalo, NY: Christian Literature Publishing Co.,；1886.；<http://www.newadvent.org/fathers/062307.htm>.}

% structure: p0001 source=h1 target=section reason=page-title

\section{十童女的筵席（第八篇）}

\Person{德克拉}\par

% structure: p0003 source=h2 target=subsection reason=relative-heading-rank; base=h2

\subsection{第一章 美多迪乌论\Quote{贞洁}一词的起源：完全属神；德性，在希腊语中}

\SourceRecord{Translated by William R. Clark.；Ante-Nicene Fathers；Vol. 6.；Edited by Alexander Roberts, James Donaldson, and A. Cleveland Coxe.；Buffalo, NY: Christian Literature Publishing Co.,；1886.；<http://www.newadvent.org/fathers/062308.htm>.}

% structure: p0001 source=h1 target=section reason=page-title

\section{\WorkTitle{童贞女的宴席（第九讲）}}

\Emph{屠西安}\par

% structure: p0003 source=h2 target=subsection reason=relative-heading-rank; base=h2

\subsection{第一章  贞洁是真帐幕的首要装饰；犹太人奉命庆祝帐棚节七天：其意义；这七天的总和不明确；无人确知世界的终结何时到来；甚至现今世界的结构已经完成。}

啊，\Person{阿瑞忒}，你是贞洁爱好者最亲爱的夸耀，我也恳求你赐予我帮助，免得我在言辞上有所欠缺，因为主题已被如此广泛而多样地探讨。因此我请求免去开场白和引言，以免我耽于适当的修饰而偏离主题：贞洁是如此光荣、可敬且著名之事。\par

天主在\BibleBook{肋未}{纪}中为真正的以色列人规定了真正帐棚节的法定礼仪，指示他们应如何庆祝和尊敬这节日；其中首要之事，是说每个人都应以贞洁装饰自己的帐棚。我将引用圣经本身的言辞，从中将毫无疑问地表明，贞洁的规诫是多么令天主喜悦和蒙祂悦纳：\BibleQuote{七月十五日，你们收了地里的出产时，应庆祝上主的节日七天：第一日为安息日，第八日也为安息日。第一日你们要拿美好树上的枝条、棕榈树枝、茂密树的枝条和溪边的柳条，在上主你们的天主面前欢乐七天。你们应每年七天庆祝上主的节日。这是你们世世代代永远的条例：你们应在七月内庆祝这节。你们应住在棚里七天；凡以色列本地人都应住在棚里，好叫你们的世代知道，当我领以色列子民出埃及时，曾使他们住在棚里：我是上主你们的天主。} \BibleRef{肋 23:39}\par

于此，犹太人如同围绕草本叶子的雄蜂，而非如同蜜蜂围绕花朵和果实，只是扑腾在圣经的字句上，完全相信这些言辞和条例是针对他们所搭的那种帐棚说的；仿佛天主喜悦于他们准备、用树木制造的那些琐碎装饰，却未觉察未来美事的丰富；然而这些事如同空气和幻影，预告了复活和我们那曾倒塌在地上的帐棚的重新支搭，最终在第七千年中，当我们再次获得不朽时，我们将在新的、不可分解的受造界中庆祝真正帐棚节的伟大盛宴，那时地上的果实已经收聚，人不再生育和被生育，而天主从创造的工程中安息。\par

因为天主在六天内创造了天和地，完成了整个世界，并在第七天休息了祂所做的一切工程，又祝福了第七天，定为圣日\BibleRef{创 2:1}，所以借着一个象征，在第七个月，当地上的果实已收聚时，我们奉命向上主庆祝节日，这表示当这个世界在第七千年结束时，当天主完成了世界，祂将在我们内喜乐。因为直到如今，万物都由祂全能的旨意和不可思议的能力所创造：大地仍出产果实，众水聚集在它们的容器中；光明仍与黑暗分开，人类的预定数目尚未完成；太阳升起统治白日，月亮统治黑夜；四足动物、走兽和爬行动物从地上生出，飞禽和水族从水中生出。那时，当指定的时期已经完成，天主停止创造这受造界时，在第七个月，即伟大的复活之日，奉命向上主庆祝我们的帐棚节，这在\BibleBook{肋未}{纪}中所说的是象征和预像；我们当仔细探究这些事，考虑赤裸的真理本身，因为祂说：\BibleQuote{智慧人听见，会增加学问；明达人会获得智谋，能明白箴言和譬喻，明白智者的言辞和他们的谜语。} \BibleRef{箴 1:5}\par

因此，让犹太人感到羞愧吧，因为他们不明白圣经的深意，以为法律和先知中除了外在事物别无所有；因为他们专注世俗之事，更看重世界的财富而非灵魂的财富。既然圣经如此划分，一部分是过去事件的写照，一部分是将来的预像，这些可怜的人却倒行逆施，把将来的预像当作已成过往之事来对待。例如祭献羔羊，他们将这奥迹仅仅视为纪念他们祖先从埃及得救，当时埃及的长子被击杀，而他们因用血涂抹家门框而得以保存。他们也不明白这预言了基督之死，借着祂的血，灵魂得以安全并被封印，将在世界焚烧时免受愤怒；而撒殚的儿子们——那些长子——将被报应天使彻底毁灭，这些天使将敬畏印在前者身上的血的印记。\par

% structure: p0009 source=h2 target=subsection reason=relative-heading-rank; base=h2

\subsection{第二章  预像、影像、真理：法律、恩宠、荣耀；人被造为不朽：死亡因毁灭性的罪而被引入。}

且容这些话作为例子，表明犹太人以何等奇妙的方式从对未来美善的希望中坠落，因为他们认为当下的事物仅仅是已成就之事的标记；他们并未察觉，预像代表图像，而图像是真理的代表。因为法律确实是图像的预像与影子，即福音的图像；但图像，也就是福音，乃是真理本身的代表。因为古人及法律向我们预言了教会的特征，而教会则代表了那将要来临的新时代的特征。因此，我们接受了基督，祂说：\Quote{我是真理} \BibleRef{若 14:16}，知道影子与预像已经终止；我们便奔向真理，宣扬它荣耀的图像。因为如今我们\Quote{所知道的只是局部}，仿佛\Quote{通过镜子观看} \BibleRef{格前 13:12}，因为那圆满的尚未临到我们；即天国与复活，那时\Quote{那局部的必要消逝} \BibleRef{格前 13:10}。因为那时我们所有的帐幕都要稳固搭建，当身体再次复活，骨头重新结合并紧密相连于血肉时。那时我们将真正向主欢庆佳节，当我们领受永恒的帐幕，不再朽坏或消散于坟墓的尘土中。如今，我们的帐幕最初固定在不可动摇的状态中，却因过犯而移动，弯曲向地，天主藉死亡终止罪恶，免得那不朽的人，作为罪人而活，且罪活在他内，遭受永罚。因此他死亡了，虽然他并非被造为可朽坏或有死的，灵魂与肉身分离，好让罪恶因死亡而灭绝，不能再存活于死者中。因此，当罪恶死亡且被消灭后，我将再次不朽地复活；我赞美天主，祂藉着死亡使祂的儿子们脱离死亡，我合法地为祂的尊荣举行佳节，以善工装饰我的帐幕，即我的肉身，正如那五位童女以五盏灯所做的那样。\par

% structure: p0011 source=h2 target=subsection reason=relative-heading-rank; base=h2

\subsection{第三章 各人应如何为将来的复活预备自己。}

在复活的第一日，我被审问是否带来了这些命令之物，是否以美德之行为装饰，是否被贞洁的枝条所荫蔽。因为你要将复活视为支搭帐幕。要将为支搭帐幕所取之物视为正义的行为。因此，我在第一日取那些规定之物，即在我站立受审的那一日，看我是否以命令之物装饰了我的帐幕；若在那一天，那些我们此时奉命预备，在那处呈献给天主的事物得以寻见。但来吧，让我们思考接下来的事。\par

\Quote{你们当取}，祂说，\Quote{在第一日美好树的枝条、棕树枝、茂密树的枝条，与溪边的柳树（以及贞洁的树），并在上主你们的天主面前欢乐。} 心未受割损的犹太人，认为最好的树木果实是香橼木，因其尺寸；他们也不以说用柏木敬拜天主为耻，即使地上的四足兽全部作为全燔祭或焚香也不足以献于祂。再者，铁石心肠的人们！若香橼在你们眼中美丽，为何不选石榴和其他树木的果实，其中包括远胜香橼的苹果呢？的确，在\WorkTitle{雅歌}中 \BibleRef{歌 4:13}，\Person{撒罗满}提到了所有这些果实，唯独对香橼缄默不语。但这欺骗了疏忽的人，因为他们不明白，生命树 \BibleRef{创 2:9} 曾由乐园所结出，如今教会又为众人所产生，即成熟的、美好的信德之果。\par

在庆节的第一日，当我们来到基督的审判台前时，我们必须带着这样的果实；因为若没有它，我们就不能与天主同庆，也不能按\Person{若望}所说 \BibleRef{默 20:6}，有分于第一次复活。因为生命树是首先被生的智慧。先知说：\Quote{她为紧握她的人是生命树，} \BibleRef{箴 3:18} \Quote{凡持守她的都是有福的。} \Quote{一棵栽在溪水旁的树，按时结果；} 即学问、爱德与明辨，适时赐予那些来到救赎之水的人。\par

凡不信基督、也不明白祂是无始之始和生命树的人，既然不能向天主展示以最美好果实装饰的帐幕，又如何能庆节呢？如何能欢乐呢？你想知道这树的美好果实吗？请深思吾主耶稣基督的言语，它们是何等甘美，超越世人。由\Person{梅瑟}而来的好果实，即法律，但并不如福音美好。因为法律是未来事物的预像和影子，但福音是真理和生命的恩宠。先知们的果实是甘美的，但不如从福音中摘取的不朽之果甘美。\par

% structure: p0016 source=h2 target=subsection reason=relative-heading-rank; base=h2

\subsection{第四章 心思清除罪恶后更清明；心思的装饰与德行次序；深而满溢的爱德；贞洁为最后的装饰；婚姻本身的运用亦当节制。}

\Quote{你们当在第一日，取美好树的枝条、棕树枝。} \BibleRef{肋 23:40} 这表示神圣纪律的操练，藉此，制服情欲的心思，通过扫除并驱赶其中的罪过而得洁净和装饰。因为我们必须清洁而装饰地前来赴宴，有如经装饰师布置，披上美德之纪律与操练。因为心思通过辛苦的操练，从那些使之昏暗的纷乱思绪中得以洁净，便迅速领悟真理；正如福音中的寡妇 \BibleRef{路 15:8}，在打扫房屋、清除灰尘之后找到了银钱；这灰尘就是那些遮蔽和笼罩心思的情欲，它们因我们的奢侈和疏忽而增长。\par

因此，谁若渴望参与那帐棚节，并愿被列入圣者之中，就当先获取信德的美果，然后取棕榈枝，即对圣经的殷勤默想与研究，之后再取那茂密多叶的爱德枝条——主在棕榈枝之后吩咐我们收取的；祂极恰当地将爱德称为浓密枝条，因为爱德全然紧密、丰盛且多结果实，没有一处稀疏或空泛，而是枝条与树干皆充实饱满。爱德正是如此，没有一处荒芜或不结果实。因为\Quote{我若将我所有的财物全施舍给穷人，并甘愿舍身被火焚烧，且具有足以移山的信德，却没有爱德，我仍不算什么。}爱德因此是众树中最浓密、最丰盛的一棵，充满且丰盈地满溢恩宠。\par

此后，祂还愿意我们取什么？柳枝；以此形象表示正义，因为按先知所言，\Quote{义人}必要\Quote{如水中之草、如溪畔之柳}那样茂盛于圣言中。最后，为加冕一切，命令将\OriginalTerm{贞洁树}{Agnos}的枝条带来装饰帐棚，因为此树以其名称即象征贞洁，使前述诸德得以美化。如今，让那些因贪恋享乐而拒绝贞洁的好色之徒离开吧。他们既未以贞洁的枝条装饰自己的帐棚，如何能与\Person{基督}一同进入节庆？这枝条是那使人神圣且蒙福的树，所有急于前往那聚会与婚宴的人都应以此束腰、遮蔽私处。因为，美丽的贞女们，请思量圣经本身及其命令：神圣之言如何将贞洁视为上述诸德与义务的冠冕，表明它在复活中是何等合宜与可欲，并且，若无贞洁，无人能获得那应许——我们这些明认守贞的人尤其应尊崇并奉献于主。那些与妻子贞洁同居的人，也拥有贞洁，他们如同在树干旁，献上低垂的枝条，承载贞洁，虽不能如我们那样触及那高大有力的树枝，甚至不能触摸它们，但他们仍真实地、虽在较低程度上，献上贞洁的枝条。然而，那些被情欲所驱策，虽未行淫，但在合法妻子的许可之事上，因未克制的情欲热力而过度拥抱的人，他们又如何庆祝节期？他们又如何欢乐？他们既未用贞洁树的枝条装饰自己的帐棚，即他们的肉身，也未听从那话：\BibleQuote{有妻子的，要像没有妻子。}\BibleRef{格前 7:29}\par

% structure: p0020 source=h2 target=subsection reason=relative-heading-rank; base=h2

\subsection{第五章 帐棚节的奥秘}

为此，我特别向那些热爱竞赛且意志坚强的人说：应当毫不迟延地尊崇贞洁，视它为最有益且最光荣的事。因为在那新的、不可分解的受造界中，凡未被贞洁枝条所装饰的，既因未按法律完成天主的命令而不得安息，又因未事先庆祝帐棚节而不得进入应许之地。因为只有庆祝了帐棚节的人，才能从那些称为帐棚的居所出发，进入圣地，直到进入天主的圣殿与圣城，迈向更大更荣耀的喜乐，一如犹太人的预表所指示。因为正如以色列人离开埃及边界后，先来到帐棚，再由此出发，进入应许之地；我们也是如此。因为我也在旅途中，从这生命的埃及出发，首先来到复活——那是真正的帐棚节——在那里支搭我的帐棚，以美德之果装饰，在复活的第一日，即审判之日，与\Person{基督}一起庆祝千年安息，称为第七日，即真正的安息日。然后，我从那里，作为\Person{耶稣}的追随者，\BibleQuote{祂已进入了诸天}\BibleRef{希 4:14}，正如以色列人在帐棚节安息后进入应许之地，我进入诸天，不再继续留在帐棚里——即我的身体不再像从前那样，而是在千年之后，从人的、可朽的形象改变为天使般的大小与美丽。在那里，我们贞女，当复活的节期达到圆满时，将从帐棚的神妙之地转向更大更好的事，上升到诸天之上的天主家中，正如圣咏作者所说：\Quote{在赞美与感谢的歌声中，在庆祝圣节的人群中。}我，阿蕾特（\Person{阿蕾特}）女士，将此礼服作为礼物献给祢，按我的能力装饰而成。\par

——欧布利奥斯（\Person{欧布利奥斯}）说：我深受感动，格里高利翁，心中思量：多米娜（\Person{多米娜}）从这些讲辞的性质来看，必是心情何等焦虑，她心中困惑，有理由害怕自己会词穷，比其他贞女讲得更乏力，因为她们在主题上说得如此有力且多样。因此，如果她明显激动，那么请你也将这补充完整；因为我好奇她作为最后一位发言者，是否有话要说。\par

——格里高利翁（\Person{格里高利翁}）说：泰奥帕特拉（\Person{泰奥帕特拉}）告诉我，欧布利奥斯，她甚为激动，但并未因词穷而困惑。因此，图西亚娜（\Person{图西亚娜}）讲完后，阿蕾特（\Person{阿蕾特}）看着她说道：来吧，我的女儿，也发表一篇讲辞，好让我们的筵席完全圆满。于是多米娜（\Person{多米娜}）脸红，迟疑许久，才勉强抬起头来，起身祈祷，转过身来，呼求智慧（\Person{智慧}）作她当下的援助。当她祈祷完毕，泰奥帕特拉说，突然勇气降临于她，一种神圣的确信占据了她，于是她说：——\par

\SourceRecord{Translated by William R. Clark.；Ante-Nicene Fathers；Vol. 6.；Edited by Alexander Roberts, James Donaldson, and A. Cleveland Coxe.；Buffalo, NY: Christian Literature Publishing Co.,；1886.；<http://www.newadvent.org/fathers/062309.htm>.}

% structure: p0001 source=h1 target=section reason=page-title

\section{《十童女的筵席》（第十篇讲道）}

\Person{多姆尼娜}\par

% structure: p0003 source=h2 target=subsection reason=relative-heading-rank; base=h2

\subsection{第一章  唯独贞洁有助于并实现最值得赞美的灵魂治理}

啊，\Person{阿勒忒}，我也省略冗长的开场白，将尽力按照我的能力进入正题，以免在那些与手头问题无关的事情上耽搁太久，说得比它们的重要性所允许的还要冗长。因为我认为审慎的一个极大部分，就是在进入主要问题之前，不要做那些仅仅取悦耳目的长篇演说，而应当立即从争论点开始。所以，我就从这里开始，因为时间已到。\par

美好的贞女们，对于德性修养，没有什么比贞洁更能使人受益了；因为唯独贞洁能够成就并带来最崇高、最美好的灵魂治理，使灵魂摆脱世上的玷污和污染，保持纯洁。因此，当基督教导我们要修持贞洁，并显明其无与伦比的美时，恶者的王国就被摧毁了——他先前曾俘虏并奴役全人类，以至于古代的人中没有一人中悦上主，反而都被错误所胜，因为法律本身不足以使人类脱离腐朽，直到贞洁接替法律，以基督的训诫治理人。如果因法律而来的正义足以使他们得救，那么最初的人们就不会如此频繁地陷入争斗和屠杀、情欲和偶像崇拜之中。实际上，他们当时被巨大而频繁的灾祸所困扰；但从基督降生成人、以贞洁武装和装饰祂的肉体之时起，那主宰不洁的残暴暴君就被废黜了，和平与信德掌权，人们不再像以前那样转向偶像崇拜。\par

% structure: p0006 source=h2 target=subsection reason=relative-heading-rank; base=h2

\subsection{第二章  解释\WorkTitle{民长纪}中树木要求立王的比喻}

但是，为了不使有些人认为我在诡辩，仅仅凭或然性推测，信口开河，我将从《旧约》中《民长纪》的预言来向你们证明，贞洁未来的统治早已被清楚地预告了。因为我们读到：\BibleQuote{有一时，树木要去立一个王管理它们，就对橄榄树说：你统治我们罢！橄榄树回答说：我岂能停止产油，用以尊崇天主和人，而去摇摆在众树之上呢？树木又对无花果树说：你来统治我们罢！无花果树回答说：我岂能停止结甜美的果实，而去摇摆在众树之上呢？树木便对葡萄树说：你来统治我们罢！葡萄树回答说：我岂能停止产新酒，使天主和人喜乐，而去摇摆在众树之上呢？于是众树对荆棘说：你来统治我们罢！荆棘对树木说：如果你们真诚地立我为王统治你们，就来投靠我的荫庇；否则，必有火从荆棘中发出，吞灭黎巴嫩的香柏木。}\par

显然，这些话并非指地上生长的树木。因为无生命的树木不能聚集商议选王，它们被深根牢牢固定在地上。这些话完全是指灵魂，它们在基督降生成人之前，过度沉溺于过犯，以恳求者的身份走近天主，请求祂的怜悯，并希望被祂的慈悲和同情所治理——圣经以橄榄为象征表现这一点，因为油对我们的身体大有裨益，能消除疲劳和疾病，并带来光明。因为所有灯火的光亮因油的滋养而增强。同样，天主的慈悲完全消解死亡，帮助人类，滋养心灵的光明。请思考，从第一个人受造直到基督，法律是否依次以这些比喻形式在圣经中被提出，而魔鬼借此欺骗了人类。圣经把无花果树比作在乐园中赐给人的命令，因为当他受骗时，他用无花果树叶子遮盖自己的赤身露体；\BibleRef{创 3:7}；把葡萄树比作在洪水时期赐给诺厄的诫命，因为他被酒所胜，受了嘲笑。\BibleRef{创 9:22}；橄榄树象征在旷野中赐给梅瑟的法律，因为当他们违反法律时，预言的恩宠——圣油——已从他们的产业中消失。最后，荆棘不无恰当地指为世界的得救而赐给宗徒们的法律；因为通过他们的训导，我们学到了贞洁，这是魔鬼唯一无法制造虚假形象的美德。因此，也有四部福音被赐下，因为天主曾四次将福音赐给人类，并通过四条法律教导他们，这些时期因果实的多样性而清楚可知。无花果树因其甘甜和丰盛，代表了人在堕落前在乐园中所拥有的快乐。确实，如我们后面将要指出的，圣神常以无花果树的果实作为美善的象征\BibleRef{耶 8:13}。葡萄树则因酒所生的喜乐，以及那些从愤怒和洪水中得救之人的欢乐，象征着从恐惧和焦虑到喜乐的转变。\BibleRef{岳 2:22}；此外，橄榄树因其产油，表示天主的慈悲，祂在洪水之后，又耐心忍受了人们转向不虔敬的行为，从而赐给他们法律，向一些人显现自己，并用油滋养那几乎熄灭的德行之光。\par

% structure: p0009 source=h2 target=subsection reason=relative-heading-rank; base=h2

\subsection{第三章  荆棘和贞洁树是贞洁的象征；四部福音，即四条教训或法律，引导人得救}

荆棘赞扬贞洁，因为荆棘与贞洁树原是同一棵树：有人称它为荆棘，有人称它为\OriginalTerm{贞洁树}{agnos}。或许因为这植物与童贞相似，故被称为荆棘和\OriginalTerm{贞洁树}{agnos}；荆棘，因其对享乐具有坚强与稳固；贞洁树，因其始终持守贞洁。因此圣经记载，厄里亚躲避依则贝耳的面，\BibleRef{列上 19:4}先来到一棵荆棘下，在那里蒙垂听，得了力量并取了食物；这表示，凡逃离情欲的煽动、逃离女人——即享乐——的人，贞洁之树便是他的避难所和荫庇，它从基督——童贞之首——来临起治理世人。因为当最初在亚当、诺厄和摩西时代颁布的法律不能赐予人救恩时，唯有福音的律法拯救了所有人。\par

这就是为何可以说无花果树未得着治理树木的王国——树木从灵性意义上指人——而无花果树象征着命令，因为人即使在堕落后仍渴望再度服从德能的统治，不愿失去乐园享乐的不朽。但人因犯罪而被弃绝，远远抛去，如同一个再不能被不朽治理、也无法承受不朽的人。犯罪后传给他的第一条信息是由诺厄宣讲的，\BibleRef{创 5:29}若他用心听从，本可免受罪恶；因为在这信息中，诺厄应许了幸福和脱离恶事的安息，只要人全力关注它，正如葡萄树应许将美酒献给殷勤栽培它的人。但这律法也没有治理人类，因为人不听从它，尽管诺厄热切宣讲。但当他们开始被水围困淹没时，他们才开始悔改，并承诺必遵守诫命。因此他们被蔑视为臣民而遭拒绝；也就是说，他们被轻蔑地告知，律法不能帮助他们；圣神回应他们，斥责他们离弃了天主命令去帮助他们、拯救他们、使他们喜乐的那些人，如诺厄和与他同在者。\Quote{连你，叛逆者，}祂说，\Quote{我来，是要帮助你们这些缺乏明智、与枯树无异的人，你们先前不信我宣讲你们应该逃离现世之事。}\par

% structure: p0012 source=h2 target=subsection reason=relative-heading-rank; base=h2

\subsection{第四章 法律无用于救恩；在荆棘形象下的最后贞洁法律。}

于是这些人在被如此拒绝于天主的眷顾之外，人类又再次陷于错谬之后，天主又藉梅瑟颁布一条法律来治理他们，召他们归向正义。但他们却决意永远告别这法律，转向了偶像崇拜。因此天主将他们交于彼此屠杀、流放和囚禁，而法律本身也犹如承认不能拯救他们。于是，他们被恶事折磨，痛苦不堪，又承诺必遵守诫命；直到天主第四次怜悯人类，派遣贞洁治理他们，圣经因此称她是荆棘。她消融享乐，并威胁说，除非众人都毫无疑问地服从她，真正来到她面前，她必用火毁灭一切，因为此后不再有其他法律或教义，只有审判和火。为此，人从此开始行正义，坚定信靠天主，并与魔鬼分离。这样，贞洁被派遣下来，因为对人最有益处和帮助。唯独她，魔鬼不能伪造仿冒以迷惑人，一如其他诫命的情况。\par

% structure: p0014 source=h2 target=subsection reason=relative-heading-rank; base=h2

\subsection{第五章 魔鬼作为万物的模仿者的恶意；两种无花果树和葡萄树。}

无花果树，如我所说，因其果实的甘甜与卓越而被视为乐园喜乐的象征；魔鬼藉其仿冒诱骗了那人，将他掳去，说服他用无花果叶遮盖身体的赤裸；也就是说，藉其摩擦激起他性欲的快乐。同样，那些从洪水中得救的人，他又用一种模仿灵性喜乐之葡萄树的饮料灌醉他们；然后再次戏弄他们，剥夺了他们的德能。我所说的话以后会更清楚。\par

仇敌凭其能力，总是模仿德能和正义的形式，并非为了真正促进其操练，而是为了欺骗和伪善。因为为了引诱那些逃离死亡的人进入死亡，他外表染上了不朽的色彩。因此他愿显得像无花果树或葡萄树，产生甘甜和喜乐，并且\BibleQuote{装作光明的天使}，\BibleRef{格后 11:14}以虔诚的外表诱骗许多人。\par

因为我们发现，在圣经中，有两种无花果树和葡萄树：\Quote{上好的无花果，非常好；极坏的无花果，非常坏。}\BibleRef{耶 24:3} \Quote{使人心愉悦的酒。}以及那毒蛇的毒汁，蝮蛇不治的毒液。\BibleRef{申 32:33}然而，自从贞洁开始掌管人以来，这欺诈就被揭露并克服了，基督，童贞女的首领，推翻了它。所以，真正的无花果树和真正的葡萄树，在贞洁的力量掌握了所有人之后，才结出果实，正如岳厄尔先知所宣告的：\Quote{土地啊，不要害怕；要欢喜快乐，因为主必要做大事。田间的野兽，不要害怕；因为旷野的草场发青，树木结果，无花果树和葡萄树都发挥力量。熙雍的子女，你们应当欢喜，因上主你们的天主而喜乐，因为他赐给了你们公义的食物。}他将先前的法律称为葡萄树和无花果树，是为属神的熙雍子女结公义之果的树木；这些树在圣言降生成人之后，当贞洁统治我们时，才结果实；而此前，因着罪恶和许多错误，它们曾抑制并毁坏了嫩芽。因为真正的葡萄树和真正的无花果树，还不能给我们提供有益于生命的滋养，那时虚假的无花果树，为欺诈而装饰多样，正在繁荣。但主枯干了虚假的枝条，即真枝条的仿制品，对那苦无花果树宣布了判决：\Quote{从今以后，你永不再结果子。}\BibleRef{玛 21:19}于是，那些真正结果实的树木才繁荣起来，结出公义的食物。\par

葡萄树，而且在多处，是指主自己，\BibleRef{若 15:1}而无花果树是指圣神，因为主\Quote{使人心愉悦}，而圣神医治他们。因此，希则克雅受命\BibleRef{列下 20:7}先用无花果饼做膏药——即圣神的果实——好使他得医治——这是按宗徒所说的——借着爱；因为他曾说：\Quote{圣神的果实是爱、喜乐、平安、忍耐、仁慈、良善、信实、温和、节制；}\BibleRef{迦 5:22}由于它们极其甘美，先知称之为无花果。米该亚也说：\Quote{他们各人要坐在自己的葡萄树和无花果树下，无人使他们害怕。}\BibleRef{米 4:4}现在，可以肯定的是，那些在圣神之下，在圣言的荫庇下，得到庇护和安息的人，将不会惊慌，也不会被那搅扰人心者吓倒。\par

% structure: p0019 source=h2 target=subsection reason=relative-heading-rank; base=h2

\subsection{第六品 匝加利亚神视的奥秘}

此外，匝加利亚表明，橄榄树影射梅瑟法律，他这样说：\Quote{那与我谈话的天使又来唤醒我，好像人从睡眠中被唤醒一样。他对我说：你看见什么？我说：我看见一个纯金的灯台，顶上有一个灯盏……旁边有两棵橄榄树，一棵在灯盏右边，一棵在左边。}\BibleRef{匝 4:1}稍后，先知询问灯台左右两边的橄榄树是什么，以及两个管子手中的两枝橄榄枝是什么，天使回答说：\Quote{这是两位受傅者，他们侍立在全地之主旁边。}这表示那两样首先的德能，它们侍奉天主，在他的居所里，通过橄榄枝，向灯芯四周供应上主的神油，使人能得见天主知识的亮光。但两棵橄榄树的两枝，就是法律和先知，围绕着那产业之分，基督和圣神是它们的作者；而在贞洁开始统治世界之前，我们自己却不能摄取这些植物的全部果实和伟大，而只能摄取它们的枝子——即法律和先知——我们从前只能培养这些，而且只是有节制地，常常让它们溜走。因为谁能接受基督或圣神，除非他先净化自己？而那从童年起就预备灵魂，使之趋向所渴望的、可喜悦的荣耀，并稳妥地带着恩宠轻松前行，从微小的劳苦中培养出伟大的希望，就是贞洁；它赋予我们的身体不朽；这是一切人所当甘心尊崇并赞美于万有之上的德行；有些人，藉着它，得以与圣言订婚，修持贞洁；另一些人，藉着它，得以脱离那诅咒：\Quote{你是尘土，仍要归于尘土。}\BibleRef{创 3:19}\par

阿瑞忒啊，这就是你所要求的关于贞洁的论述，我已尽己所能完成；我祈求，女主人，尽管它平庸而简短，愿你仁慈地从我这被选为最后发言者的人接受它。\par

\SourceRecord{Translated by William R. Clark.；Ante-Nicene Fathers；Vol. 6.；Edited by Alexander Roberts, James Donaldson, and A. Cleveland Coxe.；Buffalo, NY: Christian Literature Publishing Co.,；1886.；<http://www.newadvent.org/fathers/062310.htm>.}

% structure: p0001 source=h1 target=section reason=page-title

\section{十童女的婚筵（讲道十一）}

\Emph{\Person{阿蕾特}}\par

% structure: p0003 source=h2 target=subsection reason=relative-heading-rank; base=h2

\subsection{第一章  真实而贞洁的童女稀少；贞洁是一场竞赛；德克拉是童女之首}

特奥帕特拉叙述说，阿蕾特讲了我完全接受并赞同这一切。因为这是一件极好的事，即使你未曾如此清晰地谈论，也要认真地承担并执行这些话，不是为听众预备甜蜜的娱乐，而是为改正、提醒和节制。因为凡教导贞洁应在我的一切追求中优先被看重和拥抱的，就是正确地劝告；许多人以为他们尊敬并培养贞洁，但实话实说，真正尊敬的人很少。因为一个人若只学会抑制肉体脱离肉欲的快乐，却不能克制其余的情欲，那并不是培养贞洁；相反，他以卑鄙的私欲、以快乐交换快乐，大大地侮辱了贞洁。同样，若有人强烈抵抗感官的欲望，却被虚荣所高举，并因此能压制情欲燃烧的热度，视这一切为无物，也不能被认为是尊敬贞洁；因为当他被骄傲高举，清洗杯盘的外面，即肉体和身体，却因自负和野心伤害内心时，他就是侮辱了贞洁。再者，若有人以财富自夸，他也不会希望尊敬贞洁；他比所有人都更侮辱贞洁，因为他宁愿得到一点利益，而舍弃那在此生中被珍视的一切都无法相比的贞洁。因为一切财富与黄金，\Quote{与它相比，有如微尘。} \BibleRef{智 7:9} 同样，那过分爱自己，只急切考虑自己的益处，不顾邻居需要的人，也不是尊敬贞洁，而是侮辱贞洁。因为那从自己身上推开爱德、怜悯和仁爱的人，远逊于那些光荣地实践贞洁的人。另一方面，藉着贞洁保持童贞，却又以恶行和情欲玷污灵魂，是不对的；在这里声称纯洁和节制，在那里却因纵情于恶习而玷污它，也是不对的；同样，在这里宣称世事不给自己带来忧虑，在那里却急切地追求并关心它们，也是不对的。相反，所有的肢体都要保持完整并免受腐坏；不仅是性器官，还有那些为情欲服务的肢体。因为保持生殖器官纯洁却不保持舌头，是可笑的；或者保持舌头，却不保持视力、耳朵或手，也是可笑的；最后，保持这些纯洁却不保持心灵，以骄傲和愤怒玷污它，也是可笑的。\par

凡决心不偏离贞洁实践的人，完全有必要保持所有肢体和感官清洁并受约束，如同船上的木板，船主将它们的连接处勤勉地连结起来，免得罪恶有任何途径和通道渗入心灵。因为伟大的追求易遭大的跌倒，邪恶与真正的善对立更甚于与不善的对立。因为许多人以为抑制强烈的情欲就是贞洁，而忽略了与之相关的其他责任，结果在这事上也失败了，并将责备加于那些以正确方式追求它的人，正如你们所证明的那样，你们在一切事上是模范，在言语和行为上过着童贞的生活。现在，什么合乎童贞状态已经描述过了。\par

你们众人在我面前已在言语上充分竞赛，我宣告你们为胜利者并加冠；但是德克拉，作为你们中的首领，并且比其余众人闪耀出更明亮的光辉，要获得更大更厚的花冠。\par

% structure: p0007 source=h2 target=subsection reason=relative-heading-rank; base=h2

\subsection{第二章  德克拉得体地唱赞美诗，其余童女与她同唱；若翰洗者是贞洁的殉道者；教会是天主的净配，纯洁而童贞}

特奥帕特拉说，阿蕾特说了这些话后，命令她们全体起立，站在阿格诺斯树下，以合宜的方式向主献上感恩的赞美诗；并要德克拉开始并带领其余的人。她们站起后，她说德克拉站在阿蕾特右边的童女中间，得体地歌唱；其余的人则围成圆圈，像合唱团一样，回应她说：\Quote{我为你保守纯洁，新郎，我手持明灯前去迎接你。}\par

德克拉. 一. 从上面，童女们，唤醒死者的响声已经传来，召唤我们全体穿着白衣、手持火炬向东方去迎接新郎。醒来吧，在君王进入城门之前。\par

众贞女. 我为你保守纯洁，新郎，我手持明灯前去迎接你。\par

德克拉. 二. 逃避世人悲伤的欢乐，藐视生命奢华的享乐及其爱恋，我渴望在你赋予生命的臂膀下得到保护，并永远瞻仰你的荣美，蒙福者。\par

众贞女. 我为你保守纯洁，新郎，我手持明灯前去迎接你。\par

德克拉. 三. 为你的缘故，君王，我离开了婚姻、世人的床榻和我的金色家园，穿着无玷的衣袍前来，为能与你一同进入你有福的洞房。\par

众贞女. 我为你保守纯洁，新郎，我手持明灯前去迎接你。\par

德克拉. 四. 蒙福者，我已逃脱了蛇的无数迷人诡计，又逃脱了火焰，以及野兽毁灭凡人的袭击，我从天上等候你。\par

众贞女. 我为你保守纯洁，新郎，我手持明灯前去迎接你。\par

德克拉. 五. 主啊，因渴慕你的恩宠，我忘记了故乡。我也忘记了童女的同伴和伴侣，甚至母亲和亲人的渴望，因为基督，你是我的一切。\par

众贞女. 我为你保守纯洁，新郎，我手持明灯前去迎接你。\par

德克拉. 六. 你是生命的赐予者，基督。万福，永不沉落的光，请接受这赞美。贞女的队伍呼求你，完美的花朵、爱、喜乐、明智、智慧、圣言。\par

众贞女. 我为你保守纯洁，新郎，我手持明灯前去迎接你。\par

\Person{特克拉} 7. 敞开大门，啊，装饰华美的王后，请接纳我们进入你的内室。无玷的、荣耀得胜的新娘，散发着美丽，我们穿着如祂一般，站在基督旁边，庆祝你幸福的婚礼，啊，年轻的少女。\par

\Person{合唱} 。我为你保持纯洁，啊新郎，手持点燃的火炬，我去迎接你。\par

\Person{特克拉} 8. 站在内室外的童女们，流着苦涩的泪水，深深叹息，哀悼悲泣，因为她们的灯熄灭了，未能及时进入欢乐的内室。\BibleRef{玛 25:11}\par

\Person{合唱} 。我为你保持纯洁，啊新郎，手持点燃的火炬，我去迎接你。\par

\Person{特克拉} 9. 因为背离了神圣的生命道路，不幸的她们，忽略了为生命之路预备足够的油；手执明亮光芒已死的灯，她们从心灵深处哀叹。\par

\Person{合唱} 。我为你保持纯洁，啊新郎，手持点燃的火炬，我去迎接你。\par

\Person{特克拉} 10. 这里有盛满甜美仙酒的杯子；让我们喝吧，啊童女们，因为这是天上的饮料，是新郎为那些合法受邀参加婚礼的人准备的。\par

\Person{合唱} 。我为你保持纯洁，啊新郎，手持点燃的火炬，我去迎接你。\par

\Person{特克拉} 11. 亚伯，清楚预表你的死亡，啊受祝福者，血流如注，眼望上天，说道：\Quote{被兄弟之手残忍杀害，啊圣言，我祈求你接纳我。}\BibleRef{创 4:10}\par

\Person{合唱} 。我为你保持纯洁，啊新郎，手持点燃的火炬，我去迎接你。\par

\Person{特克拉} 12. 你勇敢的儿子若瑟，啊圣言，赢得了贞洁的最大奖赏，当我这个欲火中烧的女人强行把他拉向不洁之床时；但他不理会她，赤身逃走了，并大声呼喊：——\BibleRef{创 39:12}\par

\Person{合唱} 。我为你保持纯洁，啊新郎，手持点燃的火炬，我去迎接你。\par

\Person{特克拉} 13. 依弗太将他刚宰杀的童贞女儿祭献给天主，如同羔羊；而她，高尚地履行了祢身体的预表，啊受祝福者，勇敢地呼喊：——\par

\Person{合唱} 。我为你保持纯洁，啊新郎，手持点燃的火炬，我去迎接你。\par

\Person{特克拉} 14. 勇敢的友弟德，通过巧妙的计谋，斩下了外邦军队首领的头颅，之前她用美丽的形体引诱了他，却没有玷污自己身体的肢体，以胜利者的欢呼说道：——\par

\Person{合唱} 。我为你保持纯洁，啊新郎，手持点燃的火炬，我去迎接你。\par

\Person{特克拉} 15. 看到苏撒纳的绝美姿容，两位长老，欲火中烧，说道：亲爱的女士，我们来是想与你秘密交合；但她颤抖着呼喊：——\par

\Person{合唱} 。我为你保持纯洁，啊新郎，手持点燃的火炬，我去迎接你。\par

\Person{特克拉} 16. \Quote{我宁愿死，也不愿将我的婚礼出卖给你们这些疯癫于女色的人，并因此遭受天主永恒的公义，在烈火惩罚中。现在，啊基督，救我脱离这些邪恶。}\par

\Person{合唱} 。我为你保持纯洁，啊新郎，手持点燃的火炬，我去迎接你。\par

\Person{特克拉} 17. 你的前驱，在流溢的洗濯水中洗濯众多人群，因他的贞洁，被一个恶人无理地领去屠杀；当他用生命之血染红尘土时，他向祢呼喊，啊受祝福者：——\par

\Person{合唱} 。我为你保持纯洁，啊新郎，手持点燃的火炬，我去迎接你。\par

\Person{特克拉} 18. 你生命的母亲，那无玷的恩宠和未受玷污的童贞女，在没有男人参与的情况下怀胎，通过无玷的受孕，因此被怀疑背叛了婚床；她，啊受祝福者，怀孕时这样说道：——\BibleRef{玛 1:18}\par

\Person{合唱} 。我为你保持纯洁，啊新郎，手持点燃的火炬，我去迎接你。\par

\Person{特克拉} 19. 渴望看到你的婚礼之日，啊受祝福者，有多少天使，啊君王，祢从天上召来，带着最好的礼物来到祢面前，身穿无瑕的袍服：——\par

\Person{合唱} 。我为你保持纯洁，啊新郎，手持点燃的火炬，我去迎接你。\par

\Person{特克拉} 20. 在颂歌中，啊天主的受祝福的新娘，我们新娘的侍从尊崇你，啊无玷的童贞教会，雪白身形，黑发，贞洁，无瑕，可爱。\par

\Person{合唱} 。我为你保持纯洁，啊新郎，手持点燃的火炬，我去迎接你。\par

\Person{特克拉} 21. 腐败已经逃离，疾病的泪痛苦楚；死亡已被除去，所有愚昧消失，消耗心灵的忧伤不再存在；因为\OriginalTerm{天主基督}{God-Christ}的恩宠再次突然照耀在凡人身上。\par

\Person{合唱} 。我为你保持纯洁，啊新郎，手持点燃的火炬，我去迎接你。\par

\Person{特克拉} 22. 乐园不再与凡人隔绝，因为出于神命，他不再像从前那样居住在那里，被逐出那里，当时他尚未腐败，也没有因蛇的各种诡计而恐惧，啊受祝福者。\par

\Person{合唱} 。我为你保持纯洁，啊新郎，手持点燃的火炬，我去迎接你。\par

\Person{特克拉} 23. 唱着新歌，如今童女们伴随你走向诸天，啊王后，全都明显戴着白百合冠冕，手中拿着明亮的光。\par

\Person{合唱} 。我为你保持纯洁，啊新郎，手持点燃的火炬，我去迎接你。\par

\Person{特克拉} 24. 啊受祝福者，你从无始就居住在无玷的天庭座位，你以永恒的大能治理万物，啊圣父，与你的圣子一起，我们在这里，也请接纳我们进入生命之门。\par

\Person{合唱} 。我为你保持纯洁，啊新郎，手持点燃的火炬，我去迎接你。\par

% structure: p0057 source=h2 target=subsection reason=relative-heading-rank; base=h2

\subsection{第三章 哪些更好：节制的，还是那些喜爱生活安宁的？争战贞洁的危险：安宁的幸福；纯洁而安宁的心灵是天主：他们将看见天主；美德由试探所磨练。}

\Person{欧布利奥斯} 。当之无愧，啊格里戈里昂，特克拉赢得了首席奖赏。\par

\Person{格里戈里昂} 。确实当之无愧。\par

\Person{欧布利奥斯} 。但那位陌生人\Person{特米西亚克}呢？告诉我，她不是在门外偷听吗？我怀疑她听到这个盛宴能否保持沉默，是否不会立刻像鸟飞向食物一样，听那些所说的事？\par

\Person{格里戈里昂} 。据报告，当默多狄奥斯向\Person{阿瑞忒}询问这些事情时，她在场。但拥有像\Person{阿瑞忒}这样一位女主人和向导，即美德，是一件既美好又幸福的事。\par

\Person{Euboulios}。但是，\Person{Gregorion}，我们说哪一类人更好：是那些没有情欲而驾驭贪欲的人，还是那些在贪欲的攻击下仍保持纯洁的人？\par

\Person{Gregorion}。就我而言，我认为那些没有情欲的人是更好的，因为他们的心思纯净，完全不受玷污，在各方面都没有犯罪。\par

\Person{Euboulios}。好吧，我以贞洁起誓，你说得明智，\Person{Gregorion}。但若我反驳你的话，使你受阻碍，那是为了让我更好地学习，也为了使将来没有人能驳斥我。\par

\Person{Gregorion}。你尽管反驳我，我允许你这样做。因为，\Person{Euboulios}，我认为我足够有知识能教导你：那没有贪欲的人胜过有贪欲的人。若我不能，那就没有人能说服你了。\par

\Person{Euboulios}。天哪！我很高兴你如此宽宏地回答我，并显示出你在智慧方面是多么富有。\par

\Person{Gregorion}。你似乎只是个多嘴的人，\Person{Euboulios}。\par

\Person{Euboulios}。何以见得？\par

\Person{Gregorion}。因为你发问与其说是为了真理，不如说是为了消遣。\par

\Person{Euboulios}。请说好话吧，我的朋友；因为我非常钦佩你的智慧和声誉。我这样说是因为，关于许多智者常常互相争论的事情，你不仅说自己明白，还夸口能教导别人。\par

\Person{Gregorion}。现在请老实告诉我，你是否难以接受这个观点：那没有贪欲的人胜过那有贪欲却克制自己的人？还是你在开玩笑？\par

\Person{Euboulios}。当我告诉你我不知道时，怎么会是开玩笑呢？不过，请告诉我，最智慧的女士，那无贪欲而贞洁的人，比那有贪欲而贞洁地生活的人，好在何处？\par

\Person{Gregorion}。因为，首先，他们的灵魂本身是纯洁的，圣神常驻其中，因为它不被幻想和不受约束的思想所分心扰乱，以致玷污心灵。他们在各方面都不受情欲的侵扰，无论肉体还是心灵，都享受着情欲的安宁。而那些从外部通过视觉感官被诱惑，带着幻想，让情欲如溪流般涌入心中的人，即使他们自以为在竭力与快乐抗争并作战，也常常在心中被征服，并不少沾染污秽。\par

\Person{Euboulios}。那么我们要说，那些平静生活、不受情欲困扰的人是纯洁的吗？\par

\Person{Gregorion}。当然。因为这些人\BibleRef{玛 5:8}就是天主在真福中所称为神的人；他们毫无疑虑地相信祂。祂说他们必能坦然瞻仰天主，因为他们没有带来任何使灵魂之眼昏暗或混乱、以致不能看见天主的东西；相反，一切世俗的欲望都被除去，不仅如我所说，他们使肉体保持不与肉体结合，而且连心灵——圣神尤其在此休息和居住，如同在圣殿中——也不让任何不洁的思想进入。\par

\Person{Euboulios}。且慢；因为我认为由此我们更能顺利地发现什么才是真正最好的。告诉我，你称某人为好舵手吗？\par

\Person{Gregorion}。我当然称。\par

\Person{Euboulios}。是在巨大而令人困惑的风暴中拯救船只的人，还是在风平浪静中这样做的人？\par

\Person{Gregorion}。是在巨大而令人困惑的风暴中那样做的人。\par

\Person{Euboulios}。那么我们岂不说，那被情欲的汹涌波涛淹没，却不因此疲倦或沮丧，反而高尚地将她的小船——即肉体——驶入贞洁之港的灵魂，比在风平浪静中航行的人更好、更可敬？\par

\Person{Gregorion}。我们将这样说。\par

\Person{Euboulios}。因为预备好抵挡恶神之风的侵袭，而不被抛弃或击败，反而将一切交托给基督，并奋力与快乐抗争，这比轻松地过着贞女生活的人赢得更大的称赞。\par

\Person{Gregorion}。似乎如此。\par

\Person{Euboulios}。那么主是怎么说的？祂不是似乎表明，那虽有贪欲却保持节制的人，胜过那没有贪欲而度贞洁生活的人吗？\par

\Person{Gregorion}。祂在哪里这样说？\par

\Person{Euboulios}。在祂将智者比作根基稳固的房子，并宣告它因不能被雨水、洪水和狂风所摧毁而不可动摇时；这些风暴似是指情欲，而灵魂在贞洁上不可动摇和坚定不移的坚定似是指磐石。\par

\Person{Gregorion}。你说的似乎是真话。\par

\Person{Euboulios}。那么关于医生，你怎么说？你不是称那在大病中经受考验并治愈许多病人的医生为最好吗？\par

\Person{Gregorion}。我这样称。\par

\Person{Euboulios}。但那从未实践过、从未手中有过病人的医生，岂不是仍然处处不如？\par

\Person{Gregorion}。是的。\par

\Person{Euboulios}。那么我们可以肯定地说，一个被有贪欲的身体所包含的灵魂，以节制的药物平息因情欲之热而产生的失调，其治愈的桂冠胜过那有幸正确治理无贪欲之身体的灵魂。\par

\Person{Gregorion}。必须承认。\par

\Person{Euboulios}。摔跤又如何呢？更好的摔跤手是那个有许多强大对手并不断争斗而不被击败的人，还是那个没有对手的人？\par

\Person{Gregorion}。显然是那个摔跤的人。\par

\Person{Euboulios}。在摔跤中，不是参加比赛的运动员更有经验吗？\par

\Person{Gregorion}。必须承认。\par

\Person{Euboulios}． 因此，显然那灵魂与情欲冲动争战，不被其压倒，反而后退并摆阵对抗它的人，比那没有情欲的人显得更强。\par

\Person{Gregorion}． 真的。\par

\Person{Euboulios}． 那么怎样？\Person{Gregorion}，难道你不认为，勇敢地对抗卑劣欲望的攻击有更多的勇气吗？\par

\Person{Gregorion}． 是的，的确。\par

\Person{Euboulios}． 这勇气岂不就是美德的力量吗？\par

\Person{Gregorion}． 显然如此。\par

\Person{Euboulios}． 因此，如果忍耐是美德的力量，那么被情欲困扰却仍坚持对抗的灵魂，岂不比那不受困扰的灵魂更强？\par

\Person{Gregorion}． 是的。\par

\Person{Euboulios}． 如果更强，那就更好？\par

\Person{Gregorion}． 确实。\par

\Person{Euboulios}． 因此，从所说的看来，那有情欲并实行节制的灵魂，比那没有情欲并实行节制的灵魂更好。\par

\Person{Gregorion}． 你说得对，我更渴望与你更充分地谈论这些事。因此，如果你愿意，明天我会再来听关于它们的事。但现在，如你所见，是该照料外在人的时候了。\par

\SourceRecord{Translated by William R. Clark.；Ante-Nicene Fathers；Vol. 6.；Edited by Alexander Roberts, James Donaldson, and A. Cleveland Coxe.；Buffalo, NY: Christian Literature Publishing Co.,；1886.；<http://www.newadvent.org/fathers/062311.htm>.}
